% !TeX program = xelatex
% ============================================================
%  SY02-16 《算术与几何朗兰兹纲领》
%  底本：Xinwen Zhu, "Arithmetic and Geometric Langlands Program",
%        arXiv:2504.07502v1 [math.NT], 26 页（ICM 2022 报告综述）
%  编译：xelatex -> xelatex（无索引、无参考文献 bibtex，参考文献为 thebibliography 环境）
% ============================================================
\documentclass[UTF8, a4paper, 12pt, fontset=fandol, scheme=chinese]{ctexart}
% ============================================================
%  SY02-16 《算术与几何朗兰兹纲领》译稿导言区
%  底本：Xinwen Zhu, Arithmetic and Geometric Langlands Program,
%        arXiv:2504.07502（ICM 2022 综述，26 页）
%  体例：中文数学著作；定理类编号「节.小节.序号」，公式编号「节.序号」
% ============================================================

% ---------- 页面与字体 ----------
\usepackage{geometry}
\geometry{a4paper, left=2.6cm, right=2.6cm, top=2.8cm, bottom=2.6cm}
\linespread{1.22}
\setlength{\parskip}{0.25em}
\setlength{\emergencystretch}{1.5em}
\sloppy

% ---------- 数学 ----------
\usepackage{amsmath}
\usepackage{amsthm}
\usepackage{amssymb}
\usepackage{amsfonts}
\usepackage{mathtools}
\usepackage{mathrsfs}
\usepackage{stmaryrd}
\usepackage{bm}
\allowdisplaybreaks[2]

% ---------- 图表 ----------
\usepackage{graphicx}
\usepackage{xcolor}
\usepackage{tikz}
\usetikzlibrary{cd, arrows.meta, positioning, calc}
\usepackage[most]{tcolorbox}
\usepackage{enumitem}
\setlist{itemsep=2pt, topsep=3pt, parsep=0pt}

% ---------- 超链接 ----------
\usepackage[hyperfootnotes=false, colorlinks=true, linkcolor=blue!45!black, citecolor=green!35!black,
            urlcolor=teal!60!black, bookmarksnumbered=true]{hyperref}
\usepackage{bookmark}

% ---------- 中文微调 ----------
\xeCJKsetup{CJKecglue={\hskip 0.24em plus 0.05em minus 0.05em}}
\ctexset{
  section    = {name={,}, number=\arabic{section},
                aftername={.\enspace}, format=\Large\bfseries\centering},
  subsection = {number=\arabic{section}.\arabic{subsection},
                aftername={.\enspace}, format=\large\bfseries},
  subsubsection = {number=\arabic{section}.\arabic{subsection}.\arabic{subsubsection},
                aftername={.\enspace}, format=\normalsize\bfseries}
}

% ---------- 定理环境：共用一个计数器，按小节重置（复现原文 1.1.1 体例） ----------
\theoremstyle{plain}
\newtheorem{thm}{定理}[subsection]
\newtheorem{prop}[thm]{命题}
\newtheorem{lem}[thm]{引理}
\newtheorem{sublem}[thm]{子引理}
\newtheorem{cor}[thm]{推论}
\newtheorem{conj}[thm]{猜想}
\newtheorem{notation}[thm]{记号}
\newtheorem{construction}[thm]{构造}
\newtheorem{hypothesis}[thm]{假设}
\newtheorem{question}[thm]{问题}
\newtheorem{caution}[thm]{注意}
\newtheorem{exercise}[thm]{练习}
\theoremstyle{definition}
\newtheorem{defn}[thm]{定义}
\newtheorem{dfn}[thm]{定义}
\newtheorem{example}[thm]{例}
\newtheorem{ex}[thm]{例}
\newtheorem{convention}[thm]{约定}
\theoremstyle{remark}
\newtheorem{remark}[thm]{注}
\newtheorem{rmk}[thm]{注}

% 公式按节编号（复现原文 (1.1)、(2.3) 体例）
\numberwithin{equation}{section}
\renewcommand{\proofname}{证明}
\renewcommand{\qedsymbol}{$\square$}
\renewcommand{\refname}{参考文献}
\renewcommand{\contentsname}{目\quad 录}
\renewcommand{\indexname}{索引}

% ---------- 译注标记 ----------
% 行内脚注式译注：\zy{...}，不占用作者脚注计数器；
% 正文上标「译n」，页脚条目以「译注 n：」起头
% 译注用独立脚注序列（manyfoot），不占用作者脚注编号，标记形如「译1」
\usepackage{manyfoot}
\DeclareNewFootnote[plain]{trans}
\renewcommand{\thefootnotetrans}{译\arabic{footnotetrans}}
\newcommand{\zy}[1]{\footnotetrans{#1}}
% 段间译注框
\newtcolorbox{yizhu}[1][]{
  enhanced, breakable, colback=orange!4!white, colframe=orange!55!black,
  boxrule=0.5pt, left=6pt, right=6pt, top=4pt, bottom=4pt,
  fonttitle=\bfseries\small, title=#1,
  before skip=6pt, after skip=6pt
}

% ---------- 常用记号（原文宏的移植） ----------
% 由 工程文件/scripts/port_macros.py 从原文 preamble 自动移植，勿手改。
\newcommand{\fraka}{{\mathfrak a}}
\newcommand{\frakb}{{\mathfrak b}}
\newcommand{\frakc}{{\mathfrak c}}
\newcommand{\frakd}{{\mathfrak d}}
\newcommand{\frake}{{\mathfrak e}}
\newcommand{\frakf}{{\mathfrak f}}
\newcommand{\frakg}{{\mathfrak g}}
\newcommand{\frakh}{{\mathfrak h}}
\newcommand{\fraki}{{\mathfrak i}}
\newcommand{\frakj}{{\mathfrak j}}
\newcommand{\frakk}{{\mathfrak k}}
\newcommand{\frakl}{{\mathfrak l}}
\newcommand{\frakm}{{\mathfrak m}}
\newcommand{\frakn}{{\mathfrak n}}
\newcommand{\frako}{{\mathfrak o}}
\newcommand{\frakp}{{\mathfrak p}}
\newcommand{\frakq}{{\mathfrak q}}
\newcommand{\frakr}{{\mathfrak r}}
\newcommand{\fraks}{{\mathfrak s}}
\newcommand{\frakt}{{\mathfrak t}}
\newcommand{\fraku}{{\mathfrak u}}
\newcommand{\frakv}{{\mathfrak v}}
\newcommand{\frakw}{{\mathfrak w}}
\newcommand{\frakx}{{\mathfrak x}}
\newcommand{\fraky}{{\mathfrak y}}
\newcommand{\frakz}{{\mathfrak z}}
\newcommand{\frakA}{{\mathfrak A}}
\newcommand{\frakB}{{\mathfrak B}}
\newcommand{\frakC}{{\mathfrak C}}
\newcommand{\frakD}{{\mathfrak D}}
\newcommand{\frakE}{{\mathfrak E}}
\newcommand{\frakF}{{\mathfrak F}}
\newcommand{\frakG}{{\mathfrak G}}
\newcommand{\frakH}{{\mathfrak H}}
\newcommand{\frakI}{{\mathfrak I}}
\newcommand{\frakJ}{{\mathfrak J}}
\newcommand{\frakK}{{\mathfrak K}}
\newcommand{\frakL}{{\mathfrak L}}
\newcommand{\frakM}{{\mathfrak M}}
\newcommand{\frakN}{{\mathfrak N}}
\newcommand{\frakO}{{\mathfrak O}}
\newcommand{\frakP}{{\mathfrak P}}
\newcommand{\frakQ}{{\mathfrak Q}}
\newcommand{\frakR}{{\mathfrak R}}
\newcommand{\frakS}{{\mathfrak S}}
\newcommand{\frakT}{{\mathfrak T}}
\newcommand{\frakU}{{\mathfrak U}}
\newcommand{\frakV}{{\mathfrak V}}
\newcommand{\frakW}{{\mathfrak W}}
\newcommand{\frakX}{{\mathfrak X}}
\newcommand{\frakY}{{\mathfrak Y}}
\newcommand{\frakZ}{{\mathfrak Z}}
\newcommand{\bfa}{{\mathbf a}}
\newcommand{\bfb}{{\mathbf b}}
\newcommand{\bfc}{{\mathbf c}}
\newcommand{\bfd}{{\mathbf d}}
\newcommand{\bfe}{{\mathbf e}}
\newcommand{\bff}{{\mathbf f}}
\newcommand{\bfg}{{\mathbf g}}
\newcommand{\bfh}{{\mathbf h}}
\newcommand{\bfi}{{\mathbf i}}
\newcommand{\bfj}{{\mathbf j}}
\newcommand{\bfk}{{\mathbf k}}
\newcommand{\bfl}{{\mathbf l}}
\newcommand{\bfm}{{\mathbf m}}
\newcommand{\bfn}{{\mathbf n}}
\newcommand{\bfo}{{\mathbf o}}
\newcommand{\bfp}{{\mathbf p}}
\newcommand{\bfq}{{\mathbf q}}
\newcommand{\bfr}{{\mathbf r}}
\newcommand{\bfs}{{\mathbf s}}
\newcommand{\bft}{{\mathbf t}}
\newcommand{\bfu}{{\mathbf u}}
\newcommand{\bfv}{{\mathbf v}}
\newcommand{\bfw}{{\mathbf w}}
\newcommand{\bfx}{{\mathbf x}}
\newcommand{\bfy}{{\mathbf y}}
\newcommand{\bfz}{{\mathbf z}}
\newcommand{\bA}{{\mathbb A}}
\newcommand{\bB}{{\mathbb B}}
\newcommand{\bC}{{\mathbb C}}
\newcommand{\bD}{{\mathbb D}}
\newcommand{\bE}{{\mathbb E}}
\newcommand{\bF}{{\mathbb F}}
\newcommand{\bG}{{\mathbb G}}
\newcommand{\bH}{{\mathbb H}}
\newcommand{\bI}{{\mathbb I}}
\newcommand{\bJ}{{\mathbb J}}
\newcommand{\bK}{{\mathbb K}}
\newcommand{\bL}{{\mathbb L}}
\newcommand{\bM}{{\mathbb M}}
\newcommand{\bN}{{\mathbb N}}
\newcommand{\bO}{{\mathbb O}}
\newcommand{\bP}{{\mathbb P}}
\newcommand{\bQ}{{\mathbb Q}}
\newcommand{\bR}{{\mathbb R}}
\newcommand{\bS}{{\mathbb S}}
\newcommand{\bT}{{\mathbb T}}
\newcommand{\bU}{{\mathbb U}}
\newcommand{\bV}{{\mathbb V}}
\newcommand{\bW}{{\mathbb W}}
\newcommand{\bX}{{\mathbb X}}
\newcommand{\bY}{{\mathbb Y}}
\newcommand{\bZ}{{\mathbb Z}}
\newcommand{\mA}{{\mathcal A}}
\newcommand{\mB}{{\mathcal B}}
\newcommand{\mC}{{\mathcal C}}
\newcommand{\mD}{{\mathcal D}}
\newcommand{\mE}{{\mathcal E}}
\newcommand{\mF}{{\mathcal F}}
\newcommand{\mG}{{\mathcal G}}
\newcommand{\mH}{{\mathcal H}}
\newcommand{\mI}{{\mathcal I}}
\newcommand{\mJ}{{\mathcal J}}
\newcommand{\mK}{{\mathcal K}}
\newcommand{\mL}{{\mathcal L}}
\newcommand{\mM}{{\mathcal M}}
\newcommand{\mN}{{\mathcal N}}
\newcommand{\mO}{{\mathcal O}}
\newcommand{\mP}{{\mathcal P}}
\newcommand{\mQ}{{\mathcal Q}}
\newcommand{\mR}{{\mathcal R}}
\newcommand{\mS}{{\mathcal S}}
\newcommand{\mT}{{\mathcal T}}
\newcommand{\mU}{{\mathcal U}}
\newcommand{\mV}{{\mathcal V}}
\newcommand{\mW}{{\mathcal W}}
\newcommand{\mX}{{\mathcal X}}
\newcommand{\mY}{{\mathcal Y}}
\newcommand{\mZ}{{\mathcal Z}}
\newcommand{\scrA}{{\mathscr A}}
\newcommand{\scrB}{{\mathscr B}}
\newcommand{\scrC}{{\mathscr C}}
\newcommand{\scrD}{{\mathscr D}}
\newcommand{\scrE}{{\mathscr E}}
\newcommand{\scrF}{{\mathscr F}}
\newcommand{\scrG}{{\mathscr G}}
\newcommand{\scrH}{{\mathscr H}}
\newcommand{\scrI}{{\mathscr I}}
\newcommand{\scrJ}{{\mathscr J}}
\newcommand{\scrK}{{\mathscr K}}
\newcommand{\scrL}{{\mathscr L}}
\newcommand{\scrM}{{\mathscr M}}
\newcommand{\scrN}{{\mathscr N}}
\newcommand{\scrO}{{\mathscr O}}
\newcommand{\scrP}{{\mathscr P}}
\newcommand{\scrQ}{{\mathscr Q}}
\newcommand{\scrR}{{\mathscr R}}
\newcommand{\scrS}{{\mathscr S}}
\newcommand{\scrT}{{\mathscr T}}
\newcommand{\scrU}{{\mathscr U}}
\newcommand{\scrV}{{\mathscr V}}
\newcommand{\scrW}{{\mathscr W}}
\newcommand{\scrX}{{\mathscr X}}
\newcommand{\scrY}{{\mathscr Y}}
\newcommand{\scrZ}{{\mathscr Z}}
\newcommand{\al}{{\alpha}}
\newcommand{\be}{{\beta}}
\newcommand{\ga}{{\gamma}}
\newcommand{\de}{{\delta}}
\newcommand{\ep}{{\epsilon}}
\newcommand{\la}{{\lambda}}
\newcommand{\Ga}{{\Gamma}}
\newcommand{\La}{{\Lambda}}
\newcommand{\ve}{{\varepsilon}}
\newcommand{\Ab}{{\mathbf{Ab}}}
\newcommand{\Aff}{{\mathbf{Aff}}}
\newcommand{\Alg}{{\mathbf{Alg}}}
\newcommand{\AlgSp}{{\mathbf{AlgSp}}}
\newcommand{\Ani}{{\mathbf{Ani}}}
\newcommand{\CAlg}{{\mathbf{CAlg}}}
\newcommand{\Cat}{{\mathbf{Cat}}}
\newcommand{\Coh}{{\mathbf{Coh}}}
\newcommand{\Grp}{{\mathbf{Grp}}}
\newcommand{\Grpd}{{\mathbf{Grpd}}}
\newcommand{\infGrpd}{{\infty\on{-Groupoid}}}
\newcommand{\LCat}{{\mathbf{LCat}}}
\newcommand{\Mod}{{\mathbf{Mod}}}
\newcommand{\Mon}{{\mathbf{Mon}}}
\newcommand{\Perf}{{\mathbf{Perf}}}
\newcommand{\Perv}{{\mathbf{Perv}}}
\newcommand{\Pre}{{\mathbf{Pre}}}
\newcommand{\PStk}{{\mathbf{PStk}}}
\newcommand{\QCoh}{{\mathbf{QCoh}}}
\newcommand{\Rep}{{\mathbf{Rep}}}
\newcommand{\sCat}{{\mathbf{sCat}}}
\newcommand{\sLCat}{{\mathbf{sLCat}}}
\newcommand{\Sch}{{\mathbf{Sch}}}
\newcommand{\Sets}{\mathbf{Sets}}
\newcommand{\Shv}{{\mathbf{Shv}}}
\newcommand{\Spc}{{\mathbf{Spc}}}
\newcommand{\Stk}{{\mathbf{Stk}}}
\newcommand{\SymMon}{{\mathbf{SymMon}}}
\newcommand{\Top}{{\mathbf{Top}}}
\newcommand{\Tor}{{\mathbf{Tor}}}
\newcommand{\Vect}{{\mathbf{Vect}}}
\newcommand{\fgMod}{\Mod^{\on{f.g.}}}
\newcommand{\fgRep}{\Rep_{\on{f.g.}}}
\newcommand{\alg}{{\on{alg}}}
\newcommand{\Aut}{{\on{Aut}}}
\newcommand{\aut}{{\on{aut}}}
\newcommand{\Bcrs}{\on{B}_{\cris}}
\newcommand{\BdR}{\on{B}_{\dR}}
\newcommand{\Bst}{\on{B}_{\st}}
\newcommand{\cha}{{\on{char}}}
\newcommand{\Cha}{{\on{Char}}}
\newcommand{\coh}{{\on{coh}}}
\newcommand{\coker}{{\on{coker}}}
\newcommand{\conv}{{\on{conv}}}
\newcommand{\Corr}{{\on{Corr}}}
\newcommand{\corr}{{\on{corr}}}
\newcommand{\Der}{{\on{Der}}}
\newcommand{\der}{{\on{der}}}
\newcommand{\diag}{{\on{diag}}}
\newcommand{\End}{{\on{End}}}
\newcommand{\et}{{\on{et}}}
\newcommand{\ext}{{\on{ext}}}
\newcommand{\Ext}{{\on{Ext}}}
\newcommand{\Fil}{{\on{Fil}}}
\newcommand{\Fr}{{\on{Frob}}}
\newcommand{\Fun}{{\on{Fun}}}
\newcommand{\Gal}{{\on{Gal}}}
\newcommand{\gr}{{\on{gr}}}
\newcommand{\Gr}{{\on{Gr}}}
\newcommand{\heart}{\ensuremath\heartsuit}
\newcommand{\Hom}{{\on{Hom}}}
\newcommand{\IC}{{\on{IC}}}
\newcommand{\id}{{\on{id}}}
\newcommand{\Id}{{\on{Id}}}
\newcommand{\ind}{{\on{ind}}}
\newcommand{\Ind}{{\on{Ind}}}
\newcommand{\Int}{{\on{Int}}}
\newcommand{\inv}{{\on{Inv}}}
\newcommand{\Iso}{{\on{Isom}}}
\newcommand{\Lie}{{\on{Lie}}}
\newcommand{\Ker}{{\on{Ker}}}
\newcommand{\Map}{{\on{Map}}}
\newcommand{\mult}{{\on{mult}}}
\newcommand{\Nm}{{\on{Nm}}}
\newcommand{\Ql}{{\overline{\bQ}_\ell}}
\newcommand{\out}{\on{out}}
\newcommand{\Out}{\on{Out}}
\newcommand{\Orb}{{\on{Orb}}}
\newcommand{\ord}{{\on{ord}}}
\newcommand{\Pic}{{\on{Pic}}}
\newcommand{\pr}{{\on{pr}}}
\newcommand{\Pro}{{\on{Pro}}}
\newcommand{\proj}{{\on{proj}}}
\newcommand{\Proj}{{\on{Proj}}}
\newcommand{\pt}{{\on{pt}}}
\newcommand{\res}{{\on{res}}}
\newcommand{\Res}{{\on{Res}}}
\newcommand{\RHom}{\on{RHom}}
\newcommand{\sgn}{{\on{sgn}}}
\newcommand{\sha}{{\mbox{\textcyr{Sh}}}}
\newcommand{\Spa}{{\on{Spa}}}
\newcommand{\Spd}{{\on{Spd}}}
\newcommand{\Spec}{{\on{Spec}}}
\newcommand{\Spf}{{\on{Spf}}}
\newcommand{\std}{\on{std}}
\newcommand{\Std}{\on{Std}}
\newcommand{\sym}{\on{sym}}
\newcommand{\Sym}{\on{Sym}}
\newcommand{\tr}{{\on{tr}}}
\newcommand{\Tr}{{\on{Tr}}}
\newcommand{\xch}{\mathbb{X}^\bullet}
\newcommand{\xcoch}{\mathbb{X}_\bullet}
\newcommand{\Bun}{{\on{Bun}}}
\newcommand{\Fl}{{\mF\ell}}
\newcommand{\Loc}{{\on{Loc}}}
\newcommand{\Hk}{{\on{Hk}}}
\newcommand{\lochk}{{\Hk^{\loc}}}
\newcommand{\Sh}{{\on{Sh}}}
\newcommand{\Sht}{{\on{Sht}}}
\newcommand{\GL}{{\on{GL}}}
\newcommand{\GSp}{{\on{GSp}}}
\newcommand{\GSpin}{{\on{GSpin}}}
\newcommand{\GU}{{\on{GU}}}
\newcommand{\PGL}{{\on{PGL}}}
\newcommand{\SL}{{\on{SL}}}
\newcommand{\Sp}{{\on{Sp}}}
\newcommand{\Spin}{{\on{Spin}}}
\newcommand{\SU}{{\on{SU}}}
\newcommand{\SO}{{\on{SO}}}
\newcommand{\gl}{{\frakg\frakl}}
\newcommand{\CT}{{}^cT}
\newcommand{\CB}{{}^cB}
\newcommand{\CG}{{}^cG}
\newcommand{\LG}{{}^LG}
\newcommand{\ab}{{\on{ab}}}
\newcommand{\ad}{{\on{ad}}}
\newcommand{\Ad}{{\on{Ad}}}
\newcommand{\adm}{{\on{adm}}}
\newcommand{\aff}{{\on{aff}}}
\newcommand{\an}{{\on{an}}}
\newcommand{\Ch}{{\on{Ch}}}
\newcommand{\CH}{{\on{CH}}}
\newcommand{\cl}{{\on{cl}}}
\newcommand{\cons}{{\on{cons}}}
\newcommand{\cris}{{\on{cris}}}
\newcommand{\crys}{{\on{crys}}}
\newcommand{\cycl}{{\on{cycl}}}
\newcommand{\DCA}{\mathbf{DCA}}
\newcommand{\DerAff}{{\on{DerAff}}}
\newcommand{\disc}{{\on{disc}}}
\newcommand{\dR}{{\on{dR}}}
\newcommand{\Iw}{{\on{Iw}}}
\newcommand{\loc}{{\on{loc}}}
\newcommand{\op}{{\on{op}}}
\newcommand{\pf}{{\on{pf}}}
\newcommand{\reg}{{\on{reg}}}
\newcommand{\Tate}{{\on{Tate}}}
\newcommand{\s}{{\on{sc}}}
\newcommand{\Sat}{{\on{Sat}}}
\newcommand{\sep}{{\on{sep}}}
\newcommand{\Sph}{{\on{Sph}}}
\newcommand{\st}{{\on{st}}}
\newcommand{\St}{{\on{St}}}
\newcommand{\ta}{{\on{tame}}}
\newcommand{\un}{{\on{unip}}}
\newcommand{\ur}{{\on{ur}}}
\newcommand{\wt}{{\on{wt}}}
\newcommand{\rec}{{\on{rec}}}
\newcommand{\vect}{{\on{vec}}}
\newcommand{\kot}{{\frakB}}
\newcommand{\on}{\operatorname}


% 本文专用的少量补充
\newcommand{\Fp}{\mathbb{F}_p}


\begin{document}
\pagestyle{plain}

% ============ 扉页 ============
\begin{titlepage}
  \centering
  \vspace*{2.2cm}
  {\Huge\bfseries 算术与几何朗兰兹纲领\par}
  \vspace{0.8cm}
  {\LARGE Arithmetic and Geometric Langlands Program\par}
  \vspace{1.2cm}
  {\large 朱歆文（Xinwen Zhu）\par}
  \vspace{0.4cm}
  {\normalsize ICM 2022 报告综述\par}
  \vspace{1.6cm}
  {\large 原文：arXiv:2504.07502，26 页，2025 年\par}
  \vspace{0.5cm}
  {\large 译者：数译（DSH agent）\par}
  \vspace{0.3cm}
  {\large 版本：v1.0\qquad 日期：2026 年 11 月\par}
  \vfill
  {\small 本译稿遵循本库《学术翻译总纲》与《数学文献翻译纪律 v2.1》。\par}
\end{titlepage}

% ============ 译者说明 ============
\section*{译者说明}
\addcontentsline{toc}{section}{译者说明}

\textbf{原文性质。}\ 本篇是朱歆文（Xinwen Zhu）为 2022 年国际数学家大会（ICM）撰写的综述，刊于 ICM 2022 会议论文集，arXiv 编号 2504.07502，正文 26 页。原文不是教材，也不是专著章节，而是一份\textbf{路线图式的综述}：它逐条陈述结论（多数是定理，少数是猜想），略去证明，把"几何朗兰兹如何反过来塑造经典算术朗兰兹"这条线索讲清楚。因此本译稿严格照原文的陈述体例走：不补证明，不补背景，只在术语、记号与跳步处加译注。

\textbf{编号约定。}\ 原文定理类对象共用一个计数器并按小节重置，印刷为"节.小节.序号"（如 Theorem 1.1.4、Remark 2.1.8）；公式按节编号，印刷为"（节.序号）"（如 (1.3)、(2.2)）。本译稿完全复现这两套编号：定理类写作"定理 1.1.4"，公式写作"(1.3)"，与原文逐条对齐，正文中的交叉引用一律用 \verb|\ref| 生成，不硬编码编号。

\textbf{译注标记。}\ 需要说明的地方用两种标记：脚注式译注记作"译注 $n$"，排在当页页脚，不占用作者的脚注编号；较长的说明另设橙色方框，标题为"译注"。凡解释性文字，一律附来源（\verb|[来源：...]|）或明确标为译者推断（\verb|[译者按...]|）；来源等级按手册 A/B/I 三级执行，不使用黑名单站点。

\textbf{术语约定。}\ 术语先冻结后使用，关键几条：geometric Langlands program 几何朗兰兹纲领；Langlands dual group 朗兰兹对偶群；geometric Satake equivalence 几何 Satake 等价；perverse sheaf 反常层（沿用本库 T0836）；categorification 范畴化（沿用本库 T0140）；geometrization 几何化；decategorification 去范畴化；tamely ramified 驯顺分歧；parahoric subgroup 抛孔子群；hyperspecial 超特殊；affine Hecke category 仿射 Hecke 范畴；central sheaf 中心层；nearby cycles 邻近圈；monodromic 单值的；shtuka 保留原文 shtuka；spectral action 谱作用；Hecke eigensheaf Hecke 特征层。完整对照见 \verb|资料卡/04-术语文献勘误.md|。

\newpage
\tableofcontents
\newpage

% ============ 正文 ============

% ============================================================
%  ch00  开篇、路线图、记号与约定、致谢
%  原文对应：ICM-arXiv.tex 486--526 行
% ============================================================

经典朗兰兹纲领由 Langlands 于 1960 年代创立\cite{Langlands-to-Weil}，在表示论的框架下系统地研究互反律。粗略地说，它预测数论与表示论之间存在下述深刻联系。

\[
\begin{tikzcd}[cells={nodes={draw, inner sep=5pt}}, column sep=4.4cm, row sep=1.7cm]
  \text{Galois 表示} \arrow[r, <->, "{\scriptsize\text{互反律}}" above] \arrow[d, <->] &
  \text{自守表示} \arrow[d, <->] \\
  \text{算术数据} \arrow[r, <->, "{\scriptsize\text{Satake 同构}}" above] &
  \text{谱数据}
\end{tikzcd}
\]

这一对应的一类特例称为志村--Taniyama--Weil 猜想，它蕴含 Fermat 大定理（见 \cite{Wiles-Fermat}）。\zy{原文写作 ``Shimura-Tanniyama-Weil conjecture''，其中 Tanniyama 为 Taniyama 之误植。此处按通行拼法作 Taniyama，原样登记于 errata 表。}

几何朗兰兹纲领\cite{Laumon-Geometric-Langlands} 由 Drinfeld 与 Laumon 发起，其来源是 Drinfeld 处理函数域上 $\GL_2$ 整体朗兰兹对应的途径\cite{Drinfeld-GL2} 的推广。在几何理论中，基本研究对象从约化群 $G$ 的自守形式空间，转移到了代数曲线上 $G$-丛模空间上的层范畴。

长期以来，几何朗兰兹的发展受到经典朗兰兹中问题与技巧的启发，而另一个重要灵感来源是量子物理。其基本哲学称为范畴化／几何化。然而近些年来，几何理论反过来在经典朗兰兹纲领以及若干相关算术问题上找到了丰硕的应用。传统上，人们用 Grothendieck 的层--函数字典把几何理论中研究的层"去范畴化"，得到算术理论中研究的函数。前面提到的 Drinfeld 处理 $\GL_2$ 朗兰兹对应的工作就用了这一点。另一个著名例子是 Ngô 对基本引理的证明（\cite{Ngo-Fundamental-Lemma}）。近年来出现了从几何理论通往算术理论的另一条通道，同样借助迹构造，但性质不同，且与物理中的想法密切相关。V. Lafforgue 关于函数域上整体朗兰兹对应的工作（\cite{Lafforgue-Automorphic-to-Galois}）本质上（但并不显式地）用到了这一想法。

本文综述几何朗兰兹的（一小部分）进展，并讨论经典朗兰兹的一些近期新视角——它们由几何理论启发而来，又反过来导致若干具体算术问题的解决。下图可视为一张路线图。

\[
\begin{tikzpicture}[font=\small]
  \node[draw, inner sep=4pt, minimum width=4.8cm] (n1) {数论、算术几何、调和分析};
  \node[draw, inner sep=7pt, minimum width=4.8cm, below=11mm of n1] (n2) {\large 算术朗兰兹纲领};
  \node[draw, inner sep=7pt, minimum width=4.8cm, below=40mm of n2] (n3) {\large 几何朗兰兹纲领};
  \node[draw, inner sep=4pt, minimum width=4.8cm, below=11mm of n3] (n4) {表示论、几何、量子物理};
  \draw[<->] (n1) -- (n2);
  \draw[<->] (n3) -- (n4);
  \draw[double, -{Implies}, thick] (n2.east) .. controls +(2.5cm,0) and +(2.5cm,0) .. (n3.east)
    node[midway, right=4pt] {范畴化／几何化};
  \draw[double, -{Implies}, thick] (n3.west) .. controls +(-2.5cm,0) and +(-2.5cm,0) .. (n2.west)
    node[midway, left=4pt] {去范畴化／迹};
\end{tikzpicture}
\]

\subsection*{记号与约定}
\addcontentsline{toc}{subsection}{记号与约定}

本文通篇使用以下记号。

设 $F$ 为域，$\widetilde F/F$ 为 Galois 扩张，以 $\Ga_{\widetilde F/F}$ 记该扩张的 Galois 群。令 $\Ga_F=\Ga_{\overline F/F}$，其中 $\overline F$ 是 $F$ 的一个可分闭包。文中 $F$ 常常取局部域或整体域；此时以 $W_F$ 记 $F$ 的 Weil 群。以 $\cycl$ 记分圆特征。

设 $F$ 为域，$A$ 为 $F$ 上的乘法型群。以 $\xch(A)=\Hom(A_{\overline F},\bG_m)$ 记其特征群，以 $\xcoch(A)=\Hom(\bG_m,A_{\overline F})$ 记其余特征群。

设 $\ell$ 为素数。令 $\La$ 为 $\bF_\ell$、$\bZ_\ell$、$\bQ_\ell$，或这类环的有限（平坦）扩张。它将充当本文层理论的系数环。

\subsection*{致谢}
\addcontentsline{toc}{subsection}{致谢}

作者感谢他所有的合作者：没有他们，本综述中报告的许多工作都无从谈起。此外，他还感谢 Edward Frenkel、Dennis Gaitsgory、何旭华、Michael Rapoport、Peter Scholze 多年来在朗兰兹纲领各部分上对他的教导与讨论。本项工作部分受到 NSF（协议号 DMS-1902239）以及 Simons fellowship 的资助。

% ============================================================
%  ch01a  §1 引言 + §1.1 几何 Satake
%  原文对应：ICM-arXiv.tex 528--656 行
% ============================================================

\section{从经典朗兰兹到几何朗兰兹对应}

本节回顾几何朗兰兹理论的若干进展：其灵感来自经典理论，另一个重要来源是量子物理。基本想法是范畴化／几何化，即把集合论式的陈述替换为范畴式的相应物：
\begin{equation}\label{categorification}
\text{数}\dashrightarrow\text{向量空间}\dashrightarrow\text{范畴}\dashrightarrow\text{2-范畴}\dashrightarrow\cdots .
\end{equation}
下面用几个重要例子来说明这一过程。

\subsection{几何 Satake}\label{S: geometric Satake}

朗兰兹纲领的起点是 Satake 同构（的 Langlands 解释），其中朗兰兹对偶群神秘地登场。类似地，几何朗兰兹理论的起点是几何 Satake 等价：它是连通约化群的（球）局部 Hecke 叠上反常层范畴，与其对偶群的有限维代数表示范畴之间的张量等价。这是经典 Satake 同构的巨大推广（经由层--函数字典），并且可以说为朗兰兹对偶群（实际上是 $C$-群）何以出现在朗兰兹对应中给出了概念性的解释。

关于对偶群的记号与约定，我们遵循 \cite[Sect.~1.1]{Zhu-Integral-Satake}。设 $G$ 为域 $F$ 上的连通约化群，$(\hat G,\hat B,\hat T,\hat e)$ 为其在 $\bZ$ 上的一个带钉朗兰兹对偶群\zy{「带钉」对应原文 pinned：钉 (pinning) 指对偶群的一组结构 $(\hat B,\hat T,\hat e)$，即 Borel 子群、极大环面，以及与根向量相对应的幂幺元素 $e$；它给出同构的确定化，因而在几何 Satake 中必须一并记录。}。存在有限 Galois 扩张 $\widetilde F/F$，以及由 $\Ga_F$ 在 $G$ 的根数据上的作用诱导的单射 $\xi:\Ga_{\widetilde F/F}\hookrightarrow\Aut(\hat G,\hat B,\hat T,\hat e)$。以 $\LG=\hat G\rtimes\Ga_{\widetilde F/F}$ 记 $G$ 的通常 $L$-群，以 $\CG=\hat G\rtimes(\bG_m\times\Ga_{\widetilde F/F})$ 记 \cite{Zhu-Integral-Satake} 中定义的群，它同构于 Buzzard--Gee 引入的 $G$ 的 $C$-群。记 $d:\CG\to\bG_m\times\Ga_{\widetilde F/F}$ 为投影，其核为 $\hat G$。

设 $F$ 为非 阿基米德 局部域，$\mO$ 为其整数环，$k=\bF_q$ 为其剩余域；即 $F$ 是 $\bQ_p$ 的有限扩张，或同构于 $\bF_q(\!(\varpi)\!)$。设 $\sigma$ 为 $k$ 的\emph{几何} $q$-Frobenius 元素。假设 $G$ 可延拓为 $\mO$ 上的连通约化群（这样的 $G$ 称为非分歧的），并固定这样一个延拓，于是有 $G(\mO)\subset G(F)$，后者通常称为 $G(F)$ 的超特殊子群。以 $G(\mO)$ 的有限指数子群作为单位元的开邻域基，$G(F)$ 便成为局部紧拓扑群。经典球 Hecke 代数是 $G(F)$ 上紧支撑、$G(\mO)$-双不变的 $\bC$-值函数构成的空间，其卷积积为
\begin{equation}\label{E:conv prod alge}
(f*g)(x)=\int_{G(F)}f(y)g(y^{-1}x)dy,
\end{equation}
其中 $dy$ 是使 $G(\mO)$ 的体积为 $1$ 的 Haar 测度。注意若 $f$ 与 $g$ 都取 $\bZ$ 值，则 $f*g$ 也取 $\bZ$ 值。因此 $\bZ$ 值函数构成的子集 $H^{\cl}_{G(\mO)}$ 是一个 $\bZ$-代数\footnote{这里 $(-)^{\cl}$ 表示经典 Hecke 代数，以区别于 \eqref{E: Hecke alg isom} 中提到的导出 Hecke 代数。}。

在对偶一侧，在非分歧假设下 $\Ga_{\widetilde F/F}$ 是由 $\sigma$ 生成的有限循环群。注意 $\hat G$ 通过共轭作用在 $\CG|_{d=(q,\sigma)}$ 上，后者是 $d$ 在 $(q,\sigma)\in\bG_m\times\Ga_{\widetilde F/F}$ 处的纤维。于是经典 Satake 同构给出 $\bZ[q^{-1}]$-代数的典范同构
\begin{equation}\label{Sat isom}
\Sat^{\cl}: \bZ[q^{-1}][\CG|_{d=(q,\sigma)}]^{\hat{G}}\cong H^{\cl}_{G(\mO)}\otimes \bZ[q^{-1}].
\end{equation}

\begin{rmk}
事实上，如 \cite{Zhu-Integral-Satake} 所说明，存在定义在 $\bZ$ 上（不反转 $q$）的 Satake 同构，其中 $C$-群 $\CG$ 被替换为某个仿射幺半群，而 $\CG$ 恰为该幺半群的可逆元群。另一方面，若把基环从 $\bZ[q^{-1}]$ 扩张到 $\bZ[q^{\pm \frac{1}{2}}]$，则可把 \eqref{Sat isom} 改写成同构
\begin{equation}\label{Sat isom classical}
 \bZ[q^{\pm \frac{1}{2}}][\hat{G}\sigma]^{\hat{G}}\cong H^{\cl}_{G(\mO)}\otimes \bZ[q^{\pm \frac{1}{2}}],
\end{equation}
其中 $\hat{G}$ 通过通常的共轭作用在 $\hat{G}\sigma\subset {}^LG$ 上（讨论见 \cite{Zhu-Integral-Satake}）。这是 Satake 同构较为传统的表述，略有非典范之处，但对许多应用已经够用。
\end{rmk}

在几何理论中，人们不再把 $G(F)$ 看作拓扑群、考虑其上 $G(\mO)$-双不变紧支撑函数构成的空间，而是把 $G(F)$ 看作某个代数几何对象，研究其上 $G(\mO)$-双等变层的范畴。本节余下部分允许 $F$ 稍微更一般一些。具体地说，设 $F$ 是关于离散赋值完备的局部域，整数环为 $\mO$，剩余域 $k$ 是特征 $p>0$ 的\emph{完全}域\footnote{若 $\cha F=\cha k$（等特征情形），对 $k$ 的这一假设并不必要。此处加上它，是为了统一处理等特征与混合特征（即 $\cha F\neq \cha k$）两种情形。出于同样的理由，下文即使在等特征情形也在完全代数几何的框架下工作。}。设 $\varpi\in\mO$ 为素元。

我们在完全代数几何的框架下工作。回忆：$k$-代数 $R$ 称为完全的，指 Frobenius 自同态 $R\to R,\ r\mapsto r^p$ 是双射。以 $\Aff_k^\pf$ 记完全 $k$-代数构成的范畴。所谓完全预层（更一般地，完全预叠），指从 $\Aff_k^\pf$ 到集合范畴 $\Sets$ 的函子（更一般地，到空间构成的 $\infty$-范畴 $\Spc$ 的函子）。通常代数几何中的许多构造在此框架下依然可行。例如可以对 $\Aff_k^\pf$ 照常赋予 Zariski、étale 或 fpqc 拓扑，并谈论相应的层与叠；随后可定义完全概形、完全代数空间、完全代数叠等等，即具有某些性质的层（叠）。可以证明，这样定义的完全概形／代数空间范畴与通常意义下的完全概形／代数空间范畴等价。完全代数几何的若干基础见 \cite[Appendix A]{Zhu-Satake-Mixed}、\cite{Bhatt-Scholze-Witt} 以及 \cite[\S A.1]{Xiao-Zhu-Cycle}。

设 $R$ 为完全 $k$-代数，以 $W_\mO(R)$ 记系数在 $\mO$ 中的 $R$ 上 Witt 向量环。若 $\cha F=\cha k$，则 $W_\mO(R)\simeq R[[\varpi]]$；若 $\cha F\neq \cha k$，见 \cite[Sect.~0.5]{Zhu-Satake-Mixed}。当 $R=\overline k$ 时，记 $\mO_{\breve F}=W_\mO(\overline k)$，$\breve F=W_\mO(\overline k)[1/\varpi]$。令 $D_R=\Spec W_\mO(R)$，$D_R^*=\Spec W_\mO(R)[1/\varpi]$；可把它们看作以 $\Spec R$ 为参数的一族（打洞）圆盘。

以 $L^+G$（分别地 $LG$）记 $G$ 的\emph{射流群}（分别地\emph{环路群}）。作为 $\Aff_k^\pf$ 上的预层，
\begin{equation*}\label{E:jet and loop}
L^+G(R)=G(W_\mO(R)),\quad LG(R)=G(W_\mO(R)[1/\varpi]).
\end{equation*}
注意 $L^+G(k)=G(\mO)$，$LG(k)=G(F)$。令
\[
\Hk_G:=L^+G\backslash LG/L^+G
\]
为 $LG$ 对左、右 $L^+G$-作用的 étale 叠商，有时称为 $G$ 的（球）局部 Hecke 叠。作为完全预叠，它把 $R$ 送到三元组 $(\mE_1,\mE_2,\beta)$，其中 $\mE_1,\mE_2$ 是 $D_R$ 上的两个 $G$-挠子，而 $\beta:\mE_1|_{D_R^*}\simeq\mE_2|_{D_R^*}$ 是同构。

设 $\ell\neq p$。高阶范畴论的现代发展允许对任意预叠定义 étale $\bF_\ell$-层构成的 $\infty$-范畴（例如见 \cite{Hemo-Zhu-Unipotent}）。特别地，对 $\La=\bF_\ell,\bZ_\ell,\bQ_\ell$（或这些环的有限扩张），可以定义 $\Hk_G$ 上 $\La$-层构成的 $\infty$-范畴 $\Shv(\Hk_G,\La)$，它是 $G(F)$ 上 $G(\mO)$-双不变函数空间的范畴类比。但在不知道 $\Hk_G$ 的若干几何性质时，关于 $\Shv(\Hk_G,\La)$ 几乎无话可说。关键的几何输入是以下定理。

\begin{thm}\label{affGrass}
令 $\Gr_G:=LG/L^+G$ 为 $LG$ 对（右）$L^+G$-作用的 étale 商，它带有左 $L^+G$-作用。则 $\Gr_G$ 可写成 $L^+G$-稳定子函子的归纳极限 $\varinjlim X_i$，其中每个 $X_i$ 都是完全射影簇，且 $X_i\to X_{i+1}$ 是闭嵌入。
\end{thm}

空间 $\Gr_G$ 通常称为 $G$ 的仿射 Grassmann 簇。等特征情形见 \cite{Beauville-Laszlo, Faltings-Loop}，混合特征情形见 \cite{Zhu-Satake-Mixed, Bhatt-Scholze-Witt}；$\Gr_G$ 中闭子簇的例子见 \cite{Zhu-PCMI, Zhu-CDM}。该定理使得可以在 $\Hk_G$ 上定义可构造层与反常层的范畴，并表述几何 Satake 等价，下面就来讨论这一点。

首先，对 $L^+G$-作用取（左）商给出映射 $\Gr_G\to\Hk_G$。粗略地说，$\Hk_G$ 上的层称为反常的（分别地，可构造的），是指它拉回到 $\Gr_G$ 后来自某个 $X_i$ 上的反常（分别地，可构造）层。于是在 $\Shv(\Hk_G,\La)$ 内部有 $\Hk_G$ 上反常层与可构造层的范畴 $\Perv(\Hk_G,\La)\subset\Shv_{c}(\Hk_G,\La)$。可以把它们看作 $G(F)$ 上 $G(\mO)$-双不变紧支撑函数空间的范畴类比。此外，$\Perv(\Hk_G,\La)$ 是 Abel 范畴，当 $\La$ 是特征零的域时还是半单范畴；它称为 $G$ 的 Satake 范畴。为简单起见，下文假定 $\La$ 是域\footnote{当 $\La=\bZ_\ell$ 时相应表述略复杂一些，因为 $\mA$ 与 $\mB$ 为反常层时 \eqref{E: conv product of sheaves} 的右端未必仍是反常层。}。

卷积积 \eqref{E:conv prod alge} 也有范畴类比，即卷积图
\[
 \Hk_G\times\Hk_G\xleftarrow{\pr}  L^+G\backslash LG\times^{L^+G} LG/L^+G\xrightarrow{m} \Hk_G,
\]
而两个层 $\mA,\mB\in\Shv(\Hk_G,\La)$ 的卷积定义为
\begin{equation}\label{E: conv product of sheaves}
\mA\star\mB:= m_!\pr^*(\mA\boxtimes \mB).
\end{equation}
这一卷积积使 $\Shv(\Hk_G,\La)$ 成为幺半 $\infty$-范畴，其幺半子范畴包括 $\Perv(\Hk_G,\La)\subset \Shv_{c}(\Hk_G,\La)$。

\begin{rmk}
上述把 Satake 范畴构造为幺半范畴的方式，与更传统的做法本质等价：传统做法把 Satake 范畴定义为 $\Gr_G$ 上 $L^+G$-等变反常层的范畴（讲解见 \cite{Zhu-PCMI}）。
\end{rmk}

以 $\Coh(\bB\hat{G}_\La)^{\heart}$ 记 $\La$ 上分类叠 $\bB\hat{G}_\La$ 上凝聚层构成的 Abel 幺半范畴\footnote{在对偶群一侧，我们始终在通常代数几何的框架下工作，因此 $\bB\hat{G}$ 是通常意义下的 Artin 叠。}，它等价于 $\hat{G}$ 在有限维 $\La$-向量空间上代数表示构成的范畴。以下定理通常称为几何 Satake 等价。

\begin{thm}\label{intro:geomSat}
存在幺半 Abel 范畴之间的典范等价
\[
\Sat_G:  \Coh(\bB\hat{G}_\La)^{\heart}\cong \Perv(\Hk_G\otimes \overline k,\La).
\]
\end{thm}

几何 Satake 确实是几何朗兰兹纲领的基石之一，并已在表示论、数学物理以及（算术）代数几何中找到众多应用。当 $F=k(\!(\varpi)\!)$ 时，该定理源自 Lusztig、Ginzburg、Beilinson--Drinfeld 与 Mirković--Vilonen 的工作（参见 \cite{Lusztig-q-analog, Beilinson-Drinfeld-Hitchin, Mirkovic-Vilonen-GeoSat}）。在混合特征情形，它在 \cite{Zhu-Satake-Mixed, Yu-Integral-Satake} 中以等特征情形为输入得到证明；\cite{Fargues-Scholze-Geometrization} 则通过仿照等特征情形的策略给出证明。本节最后给出几条注记。

\begin{rmk}
\label{R: birth of the dual group}
(1) 如前所述，几何 Satake 可看作 $G$ 的朗兰兹对偶群 $\hat G$ 的概念性定义，即 Tannaka 范畴 $\Perv(\Hk_G\otimes \overline k,\La)$ 的 Tannaka 群。此外，如 \cite{Yun-Zhu-Integral-Homology, Zhu-Ramified-Sat} 所说明，$\hat G$ 典范地带有一个钉 $(\hat B, \hat T, \hat e)$。本文余下部分中，所谓 $G$ 的带钉朗兰兹对偶群 $(\hat{G},\hat{B},\hat{T},\hat{e})$，均指由几何 Satake 定义的这个四元组。

(2) 对算术应用而言，需要用对偶群一侧来理解 $\Ga_k$ 在 $\Perv(\Hk_G\otimes\overline k,\La)$ 上的作用。结果表明，这一作用由 $\Ga_k$ 在 $\hat{G}$ 上的某个作用诱导，该作用保持 $(\hat{B},\hat{T})$ 但不保持 $\hat{e}$。见 \cite{Zhu-Ramified-Sat, Zhu-PCMI}，或从 motivic 角度出发的 \cite{Zhu-CDM}。这正是群 $\CG$ 出现的原因。详细讨论见 \cite{Zhu-Ramified-Sat, Zhu-PCMI, Zhu-Integral-Satake}。

(3) 还有导出 Satake 等价 \cite{Bezrukavnikov-Finkelberg-Derived-Satake}，至少当 $\La$ 是特征零的域时，它用对偶群来描述 $\Shv_{c}(\Hk_G\otimes\overline k,\La)$。需要指出，对偶一侧出现的范畴并不是 $\bB\hat{G}_\La$ 上凝聚层的导出范畴。

(4) 事实上，对许多应用来说，重要的是几何 Satake 的族版本。设 $S$ 为（非空）有限集，则在 $D^S$ 上有局部 Hecke 叠 $\Hk_{G,D^S}$，这里 $D^S$ 是圆盘 $D=\Spec \mO$ 在 $S$ 上的自乘积。粗略地说，它分类四元组 $(\{x_s\}_{s\in S}, \mE,\mE',\beta)$，其中 $\{x_s\}_{s\in S}$ 是 $D$ 的一个 $S$-元点组，$\mE$ 与 $\mE'$ 是 $D$ 上两个 $G$-挠子，$\beta$ 是 $\mE$ 与 $\mE'$ 在 $D-\cup_{s} \{x_s\}$ 上的同构。在等特征情形，可把 $D$ 看成 $k$ 上代数曲线 $X$ 在某个 $k$-点处的形式圆盘，而 $\Hk_{G,D^S}$ 就是叠
\[
\Hk_{G, X^S}=(L^+G)_{X^S}\backslash (LG)_{X^S}/(L^+G)_{X^S}
\]
沿 $D^S\to X^S$ 的限制，其中 $(LG)_{X^S}$ 与 $(L^+G)_{X^S}$ 是 $LG$ 与 $L^+G$ 在 $X^S$ 上的族版本（精确定义见 \cite[Sect.~3.1]{Zhu-PCMI}）。在混合特征情形，叠 $\Hk_{G,D^S}$（事实上连 $D^S$ 本身）并不生活在（完全）代数几何的世界中，而生活在 Scholze 所发展的完满解析几何世界中（见 \cite{Scholze-Berkeley, Fargues-Scholze-Geometrization}）。两种情形下都可以考虑 $\Hk_{G,D^S}\otimes\overline k$ 上（ULA）反常层的范畴 $\Perv^{\on{ULA}}(\Hk_{G,D^S}\otimes\overline k,\La)$。此外，对有限集之间的映射 $S\to S'$，沿 $\Hk_{G,D^{S'}}\to\Hk_{G,D^S}$ 的限制给出函子
$\Perv^{\on{ULA}}(\Hk_{G,D^S}\otimes\overline k,\La)\to \Perv^{\on{ULA}}(\Hk_{G,D^{S'}}\otimes\overline k,\La)$\footnote{这个限制函子定义了所谓的融合积，它是几何 Satake 等价中的一个关键概念。``fusion'' 这一术语最初来自共形场论。}。细节参见上述文献。

另一方面，令 $\hat{G}^S$ 为 $\hat{G}$ 在 $\La$ 上的 $S$ 次自乘积。则对 $S\to S'$，沿 $\bB \hat{G}^{S'}\to \bB \hat{G}^S$ 的限制给出函子 $\Coh(\bB \hat{G}_{\La}^S)^\heart\to \Coh(\bB \hat{G}_{\La}^{S'})^\heart$。几何 Satake 的族版本给出函子族
\begin{equation}\label{factorizable Sat}
\Sat_S: \Coh( \bB \hat{G}_{\La}^S)^\heart\to \Perv^{\on{ULA}}(\Hk_{G,D^S}\otimes\overline k,\La),
\end{equation}
它与两侧由有限集之间的映射诱导的限制函子相容（见 \cite{Gaitsgory-deJong, Zhu-PCMI}）。

(5) 对应用而言，重要的是在不同层论框架下、在不同版本的局部 Hecke 叠上都有几何 Satake。除上面提到的之外，我们还要提到 $D$-模版本 \cite{Beilinson-Drinfeld-Hitchin}，以及算术 $D$-模版本 \cite{Xu-Zhu-Bessel}。
\end{rmk}

% ============================================================
%  ch01b  §1.2 驯顺分歧的局部几何朗兰兹对应
%  原文对应：ICM-arXiv.tex 658--748 行
%  编号：公式 (1.7)(1.8)(1.9)；注 1.2.1；定理 1.2.2；注 1.2.3；
%        定理 1.2.4；注 1.2.5；注 1.2.6
% ============================================================

\subsection{驯顺分歧的局部几何朗兰兹对应}\label{Tamely local}

先回顾经典理论。设 $F$ 是 $\bQ_p$ 的有限扩张，或同构于 $\bF_q(\!(\varpi)\!)$；为简单起见，设 $G$ 可延拓为 $\mO$ 上的连通约化群。（事实上，本小节的结果在 $F$ 的驯顺分歧扩域上分裂的拟分裂群情形下，都以适当形式成立。）此外，固定 $G$ 在 $\mO$ 上的一个钉 $(B,T,e)$。

经典局部朗兰兹纲领的目标，是用 Galois 表示来分类 $G(F)$ 的（光滑）不可约表示（系数取 $\bC$）。从这个观点看，Satake 同构 \eqref{Sat isom} 给出了 $G(F)$ 不可约非分歧表示的分类，即那些含有 $G(\mO)$-不动非零向量的表示：这类表示与 $H^{\cl}_{G(\mO)}\otimes\bC$ 的单模一一对应，而由 Satake 同构 \eqref{Sat isom}，后者又由 $\CG$ 中的半单 $\hat{G}$-共轭类参数化。（对不可约非分歧表示 $\pi$，与之对应的 $\CG$ 中的 $\hat{G}$-共轭类通常称为 $\pi$ 的 Satake 参数。）

下一类重要的不可约表示，是那些含有 Iwahori 子群不动非零向量的表示。例如，在模 $\varpi$ 约化映射 $G(\mO)\to G(k)$ 下，$B(k)\subset G(k)$ 的原像 $I$ 就是 $G(F)$ 的一个 Iwahori 子群。与非分歧情形一样，$I$-双不变的 $\bZ$ 值函数构成一个 $\bZ$-代数 $H^{\cl}_{I}$，其乘法由卷积 \eqref{E:conv prod alge} 给出（Haar 测度 $dg$ 取使 $I$ 的体积为 $1$ 者）；而 $G(F)$ 的含有非零 $I$-不动向量的不可约表示，与 $(H^{\cl}_I\otimes\bC)$-单模一一对应。

与 Satake 同构类似，Kazhdan--Lusztig 用与 $\hat{G}$ 相伴的几何对象描述了 $H^{\cl}_{I}\otimes\bC$。

以 $\hat{U}\subset\hat{B}$ 记 $\hat{B}$ 的幂幺根。自然态射 $\hat{U}/\hat{B}\to \hat{G}/\hat{G}$ 通常称为 Springer 消解。令
\[
S^{\un}_{\hat{G}}=(\hat{U}/\hat{B})\times_{\hat{G}/\hat{G}}(\hat{U}/\hat{B}),
\]
称之为 $\hat{G}$ 的（幂幺）Steinberg 叠\footnote{由于 $\hat{U}/\hat{B}\to \hat{G}/\hat{G}$ 不是平坦的，纤维积须在导出意义下取，故 $S^{\un}_{\hat{G}}$ 应理解为导出代数叠。}。在 $\bC$ 上，借助指数映射把 $\hat{U}_\bC$ 与其 Lie 代数等同，可以得到 $\hat{U}_\bC$ 上的 $\bG_{m,\bC}$-作用，从而得到 $S^{\un}_{\hat{G},\bC}$ 上的 $\bG_{m,\bC}$-作用。于是可以取商叠 $S^{\un}_{\hat{G},\bC}/\bG_{m,\bC}$。下文对 $\bC$ 上的（有限型的）Artin 叠 $X$，以 $K(X)$ 记 $X$ 上凝聚层（$\infty$-）范畴的 $K$-群。

Kazhdan--Lusztig \cite{Kazhdan-Lusztig-Iwahori} 在 $G$ 分裂且中心连通的假设下，构造了典范同构（需先取定 $q$ 的一个平方根 $\sqrt{q}$）
\begin{equation}\label{KL isomorphism}
K(S^{\un}_{\hat{G},\bC}/\bG_{m,\bC})\otimes_{K(\bB\bG_{m,\bC})}\bC\cong H^{\cl}_I\otimes\bC,
\end{equation}
其中映射 $K(\bB\bG_{m,\bC})\to \bC$ 把与 $\bG_{m,\bC}$ 的重言表示相对应的类送到 $\sqrt{q}$。此外，该同构诱导出 Bernstein 同构
\begin{equation}\label{Bernstein isomorphism}
K(\bB \hat{G}_\bC)\otimes\bC\cong Z(H^{\cl}_I\otimes\bC),
\end{equation}
其中 $Z(H^{\cl}_I\otimes\bC)$ 是 $H^{\cl}_I\otimes\bC$ 的中心，而映射 $K(\bB \hat{G}_\bC)\to K(S^{\un}_{\hat{G},\bC}/\bG_{m,\bC})$ 由自然投影 $S^{\un}_{\hat{G}}/\bG_m\to \bB\hat{G}$ 诱导。

\begin{rmk}
用涉及 $\hat{G}$ 的几何来刻画 $\bZ$-代数 $H^{\cl}_I$，将是很有意思的问题；这将推广 \cite{Zhu-Integral-Satake} 中的整系数 Satake 同构。
\end{rmk}

Kazhdan--Lusztig 同构 \eqref{KL isomorphism} 也容许范畴化，通常称为 Bezrukavnikov 等价：它给出仿射 Hecke 范畴的两种实现。与前面一样，转入几何理论时允许 $F$ 稍更一般，如 \S\ref{S: geometric Satake} 中那样；同时设 $G$ 可延拓为 $\mO$ 上的连通约化群，并固定 $G$ 在 $\mO$ 上的一个钉。

设 $L^+G\to G_k$ 为自然的“模 $\varpi$ 约化”映射，以 $\Iw\subset L^+G$ 记 $B_k\subset G_k$ 的原像。这是 $I$ 的几何类比。于是如 \S\ref{S: geometric Satake} 中讨论的非分歧情形，可以定义 Iwahori 局部 Hecke 叠 $\Hk_{\Iw}= \Iw\backslash LG/\Iw$ 与幺半范畴
$\Shv_{c}(\Hk_{\Iw}\otimes\overline k,\La)\subset\Shv(\Hk_{\Iw}\otimes\overline k,\La)$。范畴 $\Shv_{c}(\Hk_{\Iw}\otimes\bar k,\La)$ 可看作 $H^{\cl}_I$ 的范畴类比，通常称为仿射 Hecke 范畴。

回忆 $\breve F=W_\mO(\overline k)[1/\varpi]$。$F$ 的惯性群 $I_F:=\Ga_{\breve F}$ 有驯顺商 $I_F^t$，同构于 $\prod_{\ell\neq p} \bZ_\ell(1)$。

\begin{thm} \label{T: Bezrukavnikov equivalence}
对驯顺惯性群 $I_F^t$ 的任一拓扑生成元 $\tau$ 的选取，存在幺半 $\infty$-范畴之间的典范等价
\[
\on{Bez}^{\un}_G:   \Coh(S^{\un}_{\hat{G},\La})\cong \Shv_{c}(\Hk_{\Iw}\otimes\overline k,\La).
\]
\end{thm}

事实上，Bezrukavnikov 在 $F=k(\!(\varpi)\!)$ 时于 \cite{Bezrukavnikov-Affine-Hecke} 中证明了这一等价；Yun 与作者由等特征情形推出混合特征情形。一个有意思的问题是：\cite{Scholze-Berkeley, Fargues-Scholze-Geometrization} 引入的新技术能否直接给出该等价在混合特征情形的证明。（这方面的若干进展见 \cite{AGLR}。）

\begin{rmk}
再次指出，对算术应用而言，需要用对偶群一侧来描述 $\Ga_k$ 在 $\Shv_{c}(\Hk_{\Iw}\otimes\bar k,\La)$ 上的作用。讨论见 \cite{Bezrukavnikov-Affine-Hecke, Hemo-Zhu-Unipotent}。
\end{rmk}

我们说明定理 \ref{T: Bezrukavnikov equivalence} 的证明（在 $F=k(\!(\varpi)\!)$ 时）中的一个重要环节。存在 $\mO$ 上光滑仿射群概形 $\mG$（称为 Iwahori 群概形），使得 $\mG\otimes F=G$ 且 $L^+\mG=\Iw$。于是有 $D$ 上的局部 Hecke 叠 $\Hk_{\mG,D}$，它类似于 \S\ref{S: geometric Satake} 末尾讨论的 $\Hk_{G,D^S}$（此处 $S=\{1\}$）。此外 $\Hk_{\mG,D}|_{D^*}\cong \Hk_{G,D}|_{D^*}$，且 $\Hk_{\mG,D}|_0=\Hk_{\Iw}$，其中 $0\in D$ 是闭点。

取邻近圈便得到
\begin{equation}\label{E: central sheaf}
\mZ: \Coh(\bB\hat{G}_\La)^\heart\xrightarrow{\Sat_{\{1\}}} \Perv(\Hk_{\mG,D}|_{D^*_{\bar k}},\La)\xrightarrow{\Psi} \Perv(\Hk_{\Iw}\otimes \bar k,\La).
\end{equation}
这就是 Gaitsgory 的中心函子 \cite{Gaitsgory-Central-Sheaf, Zhu-Coherence-Conjecture}，可视为 \eqref{Bernstein isomorphism} 的范畴化。这一构造的动机来自研究志村簇 mod $p$ 几何时提出的 Kottwitz 猜想，相关讨论见 \S\ref{S: local model}。

定理 \ref{T: Bezrukavnikov equivalence} 可以推广到驯顺层。考虑交换图
\[
\hat{G}/\hat{G}\leftarrow \hat{B}/\hat{B}\xrightarrow{q_{\hat{B}}} \hat{T}/\hat{T},
\]
其中左态射是通常的 Grothendieck--Springer 消解。设 $\chi$ 是 $\hat{T}/\hat{T}$ 的一个 $\La$-点，其中 $\La$ 是 $\bQ_\ell$ 的有限扩张。令 $(\hat{B}/\hat{B})_\chi=q_{\hat{B}}^{-1}(\chi)$，并令
\[
S^{\chi}_{\hat{G},\La}:= (\hat{B}/\hat{B})_\chi\times_{\hat{G}/\hat{G}} (\hat{B}/\hat{B})_\chi.
\]
注意当 $\chi=1$ 时它化归为 $S^{\un}_{\hat{G},\La}$。

另一方面，（挠的）$\La$-点 $\chi\in \hat{T}/\hat{T}$ 定义了 $\Iw\otimes\overline{k}$ 上的一维特征层 $\mL_\chi$。于是可以定义 $LG_{\bar k}$ 上双 $(\Iw,\mL_\chi)$-等变可构造层的幺半范畴，记作 $\Shv_{\cons}({}_\chi(\Hk_{\Iw})_{\chi},\La)$。当 $\chi=1$ 时 $\mL_\chi$ 就是 $\Iw$ 上的平凡特征层，此即仿射 Hecke 范畴 $\Shv_c(\Hk_{\Iw}\otimes\overline k,\La)$。定理 \ref{T: Bezrukavnikov equivalence} 的如下推广由 \cite{Bezrukavnikov-Affine-Hecke} 提出，它将在与 Dhillon--Li--Yun 的合作文章 \cite{DLYZ-Endoscopy} 中证明。

\begin{thm} \label{T: Monodromic Bezrukavnikov equivalence}
设 $\cha F=\cha k$。存在典范的幺半等价
\[
\on{Bez}^{\chi}_{\hat{G}}: \Coh(S^{\chi}_{\hat{G},\La})\cong \Shv_{c}({}_\chi(\Hk_{\Iw})_{\chi},\La).
\]
\end{thm}

\begin{rmk}\label{R: int Bez}
为定理 \ref{T: Bezrukavnikov equivalence} 与定理 \ref{T: Monodromic Bezrukavnikov equivalence} 中的等价建立 $\bZ_\ell$-层版本，是重要的。
\end{rmk}

\begin{rmk}\label{2-categorical LGL}
驯顺分歧之外的局部几何朗兰兹对应尚未被完全理解，尽管某些野分歧已在具体问题中出现（例如 \cite{HNY-Kloosterman, Zhu-FG-vs-HNY}）。人们普遍相信，一般的局部几何朗兰兹应以 2-范畴式的陈述来表述：它预测，在 $LG$ 上（适当定义的）层范畴作用下，模范畴构成的 2-范畴，与局部几何朗兰兹参数叠上的范畴构成的 2-范畴等价。确切表述超出本文范围；粗略地说，它蕴含（并且多少等价于说）$LG$ 在适当选取的“水平”下的 Hecke 范畴（Morita）等价于某个形如 $X\times_{Y}X$ 的叠上凝聚层范畴，其中 $Y$ 与局部几何朗兰兹参数的模空间密切相关。
\end{rmk}

% ============================================================
%  ch01c  §1.3 整体几何朗兰兹对应
%  原文对应：ICM-arXiv.tex 750--780 行
%  编号：猜想 1.3.1
% ============================================================

\subsection{整体几何朗兰兹对应}\label{Global GLC}

如本文开头所述，（整体）几何朗兰兹纲领源自 Drinfeld 对函数域上 $\GL_2$ 朗兰兹猜想的证明。这一学科的早期发展，大多集中于对整体函数域 $F$ 的 Galois 表示（等价地，光滑代数曲线 $X$ 上的局部系统）构造 Hecke 特征层，参见 \cite{Drinfeld-GL2, Laumon-Geometric-Langlands, FGV-Geometric-Langlands-GLn}。

此后，在 \cite{Beilinson-Drinfeld-Hitchin} 的工作之后，整个纲领的范围发生了转变：Beilinson--Drinfeld 在其中给出了整体几何朗兰兹对应的一个粗略范畴式表述。这一表述随后由 Arinkin--Gaitsgory 在 \cite{Arinkin-Gaitsgory-Singular-Support} 中精确化，下面加以回顾。

设 $X$ 是 $F=\bC$ 上的光滑射影曲线。在自守一侧，令 $\mathbf{D}_c(\Bun_G)$ 为 $X$ 上主 $G$-丛的模叠 $\Bun_G$ 上凝聚 D-模构成的 $\infty$-范畴。在 Galois 一侧，令 $\Coh(\Loc_{\hat{G}})$ 为 $X$ 上 de Rham $\hat{G}$-局部系统（又称带平坦联络的主 $\hat{G}$-丛）的模叠 $\Loc_{\hat{G}}$ 上凝聚层构成的 $\infty$-范畴。

\begin{conj}\label{Conj: GGLC}
存在 $\infty$-范畴之间的典范等价
\[
\bL_G:  \Coh(\Loc_{\hat{G}})\cong \mathbf{D}_c(\Bun_G),
\]
且满足一列自然的相容性条件。
\end{conj}

我们只简述其中最重要的相容性条件，其余见 \cite{Arinkin-Gaitsgory-Singular-Support}。注意两侧都带有一族相互交换的算子的作用，这些算子由 $x\in X$ 与 $V\in \Coh(\bB\hat{G}_\bC)^\heart$ 标号。具体地说，对每一点 $x\in X$ 有求值映射 $\Loc_{\hat{G}}\to \bB\hat{G}_\bC$，于是每个 $V\in \Coh(\bB\hat{G}_\bC)^\heart$ 通过拉回给出 $\Loc_{\hat{G}}$ 上的向量丛 $\widetilde{V}_x$，后者再以张量积作用于 $\Coh(\Loc_{\hat{G}})$。另一方面，有 Hecke 算子 $H_{V,x}$ 作用于 $\mathbf{D}_c(\Bun_G)$，其做法是借助（几何 Satake \eqref{factorizable Sat} 的 D-模版本）把层 $\Sat_{\{1\}}(V)|_x$ 作卷积。等价 $\bL_G$ 应当把这些算子的作用相互对应起来。

尽管这一猜想仍远未解决，人们已知 $\Loc_{\hat{G}}$ 上完全复形的范畴 $\Perf(\Loc_{\hat{G}})$ 作用于 $\mathbf{D}_c(\Bun_G)$，该作用通常称为谱作用；在此作用下，$\widetilde{V}_x\in \Perf(\Loc_{\hat{G}})$ 对 $\mathbf{D}_c(\Bun_G)$ 的作用由 Hecke 算子 $H_{V,x}$ 给出。

如今猜想 \ref{Conj: GGLC} 有时被称为整体几何朗兰兹猜想的 de Rham 版本，因为这类猜想式等价还有其他版本，下面略作介绍。

首先，本着非交换 Hodge 理论的精神，猜想 \ref{Conj: GGLC} 应当存在一个经典极限，有时称为整体几何朗兰兹的 Dolbeault 版本。其精确表述（据作者所知）尚不清楚，但一般说来，它归结为 $G$ 与 $\hat{G}$ 的 Hitchin 纤维化的对偶（在镜像对称的意义下）；这一点已在 \cite{Donagi-Pantev-Hitchin, Chen-Zhu-Geometric-Langlands-char-p} 中“一般地”建立。通过在正特征下扭变／形变这种对偶，可以证明猜想 \ref{Conj: GGLC} 的特征 $p$ 类比（当然也只是“一般地”，见 \cite{Bezrukavnikov-Braverman-Geometric-Langlands-char-p, Chen-Zhu-Nonabelian-Hodge, Chen-Zhu-Geometric-Langlands-char-p}）。有意思的是，这个“一般”的特征 $p$ 等价可用来给出 \cite{Beilinson-Drinfeld-Hitchin} 主要结果的新证明（至少对 $G=\GL_n$，见 \cite{Bezrukavnikov-Travkin-Quantization}）。

\cite{Beilinson-Drinfeld-Hitchin} 的工作（从而整体几何朗兰兹的 de Rham 版本）深受共形场论的影响。另一方面，受拓扑场论启发，Ben-Zvi 与 Nadler \cite{Ben-Zvi-Nadler-Betti-Geometric-Langlands} 提出了猜想 \ref{Conj: GGLC} 的 Betti 版本：在自守一侧，把 $\Bun_G$ 上 $D$-模范畴替换为（$\Bun_G$ 解析化后的）$\bC$-向量空间层范畴；在 Galois 一侧，把 $\Loc_{\hat{G}}$ 替换为 Betti $\hat{G}$-局部系统的模空间（又称 $X$ 的基本群的 $\hat{G}$-值表示）。

Riemann--Hilbert 对应允许在 de Rham $\hat{G}$-局部系统与 Betti $\hat{G}$-局部系统之间过渡，但这一过渡是超越的。所以猜想 \ref{Conj: GGLC} 与其 Betti 类比并无直接联系。

最近，Arinkin 等人 \cite{AGKRRV} 提出了猜想 \ref{Conj: GGLC} 的又一个版本，它直接把 de Rham 版本与 Betti 版本联系起来，同时又包含用 $\ell$-进层表述的版本。因此它与函数域上的经典朗兰兹对应关系更密切，这将在 \S\ref{S: Global Langlands function fields} 中讨论。

% ============================================================
%  ch02a  §2 引言 + §2.1 范畴算术局部朗兰兹
%  原文对应：ICM-arXiv.tex 782--935 行
%  编号：定理 2.1.1；定理 2.1.2；猜想 2.1.3；注 2.1.4；
%        公式 (2.1)(2.2)(2.3)(2.4)；定理 2.1.5；例 2.1.6；例 2.1.7；注 2.1.8
% ============================================================

\section{从几何到经典朗兰兹纲领}

上一节讨论了范畴化与几何化的想法如何引出几何朗兰兹纲领的发展。另一方面，量子物理的想法允许我们在不同维数的流形上取值（拓扑）量子场论，从而把 \eqref{categorification} 中的箭头反向。这类想法在几何与拓扑中当然并不新鲜；但出人意料的是，它也带来了对经典朗兰兹纲领的新理解。事实上，众所周知整体域与 $3$-流形之间存在类比，在此类比下 Frobenius 对应于圆周的基本群。于是“场论在圆周上的紧化”引出范畴迹方法（例如见 \cite{AGKRRV, BCHN, Zhu-CDM}），它在几何表示论中起着越来越重要的作用。

\subsection{范畴算术局部朗兰兹}\label{S: Local Langlands}

本小节中，设 $F$ 是 $\bQ_p$ 的有限扩张，或同构于 $\bF_q((\varpi))$。经典局部朗兰兹对应寻求用 Galois 数据来分类 $G(F)$ 的光滑不可约表示。除 $G=\GL_n$ 情形（这是 \cite{LRS-GLn, Harris-Taylor-GLn} 的定理）之外，其精确表述相当复杂。然而，局部几何朗兰兹是 2-范畴式的这一信条（见注 \ref{2-categorical LGL}）提示：即便经典局部朗兰兹对应也应当、而且多半需要被范畴化。

表述范畴算术局部朗兰兹所需的第一个要素是下述结果，它由 \cite{DHKM, Fargues-Scholze-Geometrization, Zhu-Conjecture} 各自独立得到。我们采用 \cite{Zhu-Conjecture} 中的表述，其中（强）连续同态的概念参见该文。

\begin{thm}\label{Stack of local parameters}
把 $\bZ_\ell$-代数 $A$ 送到满足 $d\circ\rho=(\cycl^{-1},\pr)$ 的（强）连续同态 $\rho: W_F\to \CG(A)$ 所构成空间的预叠，可由一个（经典）概形 $\Loc_{\CG}^\Box$ 表示；该概形是若干仿射概形的无交并，每个仿射概形在 $\bZ_\ell$ 上平坦、有限型，且局部完全交。
\end{thm}

$\hat{G}$ 在 $\CG$ 上的共轭作用诱导出 $\hat{G}$ 在 $\Loc_{\CG}^\Box$ 上的作用；称商叠 $\Loc_{\CG}=\Loc_{\CG}^\Box/\hat{G}$ 为局部朗兰兹参数叠，它粗略地说把上述诸 $\rho$ 构成的群胚按 $\hat{G}$-共轭分类。

在局部朗兰兹对应的范畴版本中，Galois 一侧自然要考虑 $\Loc_{\CG}$ 上凝聚层构成的（$\infty$-）范畴 $\Coh(\Loc_{\CG})$。表示一侧，人们或许会朴素地考虑 $G(F)$ 光滑表示构成的（$\infty$-）范畴 $\Rep(G(F),\La)$。但实际上这个范畴需要扩大。

这可以从不同角度看出。事实上，许多人共有的一条普遍经验是：在经典局部朗兰兹对应中，最好把 $G$ 的表示与一族群（它们是 $G$ 的内形式的精细化）放在一起研究。这类群的族有各种提法。算术几何（即对志村簇与 Shtuka 模空间的 $p$-进与 mod $p$ 几何的研究）与几何表示论（即范畴迹构造）都提示，应当研究由 Kottwitz 集
\[
B(G) = G(\breve F)/\sim, \quad \mbox{其中 } g\sim g' \mbox{ 指存在 } h\in G(\breve F) \mbox{ 使 } g'=h^{-1}g\sigma(h),
\]
所给出的一族群 $\{J_b(F)\}_{b\in B(G)}$ 的表示范畴粘合而成的范畴。这里对 $b\in B(G)$（提升为 $G(\breve F)$ 中的元素），群 $J_b$ 是按下述方式定义的 $F$-群：对 $F$-代数 $R$，令 $J_b(R)=\{h\in G(\breve F\otimes_FR)\mid h^{-1}b\sigma(h)=b\}$。特别地，若 $b=1$ 则 $J_b=G$。一般地，存在良定义到共轭的嵌入 $(J_{b})_{\overline{F}}\to G_{\overline{F}}$，它使 $J_b$ 成为 $G$ 的某个 Levi 子群的内形式的精细化（设 $G$ 拟分裂）。

有两种方式把这想法精确化。一种属于 Fargues--Scholze \cite{Fargues-Scholze-Geometrization}：他们把 $B(G)$ 看作 Fargues--Fontaine 曲线上 $G$-丛的 $v$-叠 $\Bun_G$ 的点集，并考虑 $\Bun_G$ 上适当定义的 étale 层构成的范畴 $D_{\on{lis}}(\Bun_G,\La)$，后者确实把所有 $\Rep(J_b(F),\La)$ 粘合在一起。需要指出，这一途径依赖 Scholze 关于 diamond 与凝聚态数学的 $\ell$-进形式体系的工作。

另一途径 \cite{Xiao-Zhu-Cycle, Zhu-CDM, Zhu-Conjecture, Hemo-Zhu-Unipotent} 与范畴迹的想法密切相关：它把 $B(G)$ 看作（étale）商叠
\begin{equation*}\label{E: kot stack}
\frakB(G):=LG/\Ad_\sigma LG
\end{equation*}
的点集，其中 $\Ad_\sigma$ 表示由
\[
\Ad_\sigma: LG\times LG\to LG,\quad  (h,g)\mapsto hg\sigma(h)^{-1}
\]
给出的 Frobenius 扭共轭。于是有前面提到的 $\La$-层范畴 $\Shv(\kot(G)\otimes \bar{k},\La)$。尽管 $\kot(G)$ 在传统代数几何中是病态对象，关于它的几何仍可以说上几句，而范畴 $\Shv(\frakB(G)\otimes \bar{k},\La)$ 相当合理。此外，还可以定义 $\frakB(G)\otimes \bar{k}$ 上可构造层构成的范畴 $\Shv_c(\frakB(G)\otimes \bar{k},\La)$，下面略作说明，细致讨论见 \cite{Hemo-Zhu-Unipotent}。

对每个 $k$ 上的代数闭域 $\Omega$，$\kot(G)$ 的 $\Omega$-点群胚就是 $\Omega$ 上带 $G$-结构的 $F$-同晶体的群胚，其同构类集合同于 Kottwitz 集 $B(G)$。然而 $\kot(G)$ 并不只是其各点的无交并；它带有一个分层，称为 Newton 分层，以 $B(G)$ 标号。具体地说，$B(G)$ 上有自然的偏序，粗略地讲，对每个 $b\in B(G)$，对应于 $b'\leq b$ 的那些 $\Omega$-点构成闭子叠 $i_{\leq b}: \kot(G)_{\leq b}\subset \kot(G)\otimes \bar k$，而对应于 $b$ 的那些 $\Omega$-点构成开子叠 $j_{b}: \kot(G)_b\subset \kot(G)_{\leq b}$。

特别地，$B(G)$ 中的基本元（即关于偏序 $\leq$ 的极小元）给出闭层。还要指出，若 $b$ 是基本元，则 $J_b$ 是 $G$ 的某个内形式的精细化，通常称为 $G$ 的扩张纯内形式。

本小节余下部分把 $\kot(G)\otimes \bar k$ 简记为 $\kot(G)$，并记 $i_b=i_{\leq b}j_b: \kot(G)_b\hookrightarrow \kot(G)$ 为局部闭嵌入。对 $b$，以 $\fgRep(J_b(F),\La)$ 记 $\Rep(J_b(F),\La)$ 中由开紧子群 $K\subset J_b(F)$ 的平凡表示紧诱导出的表示
\[
\delta_{K,\La}:=c\on{-ind}_K^{J_b(F)}(\La)
\]
在有限余极限与收缩下生成的完全子范畴。下述定理取自 \cite{Hemo-Zhu-Unipotent}，它概括了 $\Shv_c(\kot(G),\La)$ 的若干性质。

\begin{thm}\label{T: Local-Langlands-Category}
\begin{enumerate}
    \item $\Shv(\kot(G),\La)$ 中的对象可构造，当且仅当它对每个 $\frakB(G)_b$ 的 $!$-限制可构造，且对几乎所有 $b$ 为零。若 $\La$ 是特征零的域，则 $\Shv_{c}(\frakB(G),\La)$ 由 $\Shv(\kot(G),\La)$ 的紧对象组成。
    \item 对每个 $b\in B(G)$，存在典范等价 $\Shv_c(\frakB(G)_b,\La)\cong \fgRep(J_b(F),\La)$。存在完全忠实嵌入 $i_{b,*}, i_{b,!}:  \Shv_c(\frakB(G)_b,\La)\to \Shv_c(\kot(G),\La)$（当 $b$ 为基本元时两者相同），它们给出 $\Shv_c(\frakB(G),\La)$ 关于 $\{\Shv_c(\frakB(G)_b,\La)\}_b$ 的半正交分解。
    \item 存在自对偶函子 $\bD^{\coh}\colon \Shv_c(\frakB(G),\La)\simeq \Shv_c(\frakB(G),\La)^{\vee}$，其构造方式是把各个 $\fgRep(J_b(F),\La)$ 上的上同调对偶（Bernstein--Zelevinsky 意义下）粘合起来。
    \item 把各个 $\fgRep(J_b(F),\La)$ 上（平移后的）$t$-结构粘合起来，得到自然的反常 $t$-结构；当 $\La$ 是域时，它被 $\bD^{\coh}$ 保持。
\end{enumerate}
\end{thm}

下述算术局部朗兰兹猜想的范畴形式（\cite[Sect. 4.6]{Zhu-Conjecture}）受到 \S\ref{Global GLC} 中讨论的整体几何朗兰兹猜想的启发。

\begin{conj}\label{C: LLC}
设 $G$ 在 $F$ 上拟分裂，并带有钉 $(B,T,e)$；固定非平凡加性特征 $\psi: F\to \bZ_\ell[\mu_{p^\infty}]^\times$。则存在范畴之间的典范等价
\[
\bL_{G}:   \Coh(\Loc_{\CG}\otimes\La)\cong \Shv_{c}(\frakB(G),\La).
\]
\end{conj}

\begin{rmk}
(1) 该猜想有一个密切相关的版本：把 $\Shv_{c}(\frakB(G),\La)$ 换成 $\Shv(\frakB(G),\La)$，把 $\Coh(\Loc_{\CG}\otimes\La)$ 换成其 ind-完备化（并加上某些支撑条件）（见 \cite[Sect. 4.6]{Zhu-Conjecture}）。Fargues--Scholze \cite{Fargues-Scholze-Geometrization} 提出了一个与此版本平行的猜想，其中范畴 $\Shv(\frakB(G),\La)$ 被换成了前面提到的 $D_{\on{lis}}(\Bun_G,\La)$。

(2) \cite{Zhu-Conjecture} 还说明了在 $\bQ$ 上建立此类等价的一个 motivic 期望。
\end{rmk}

该猜想的一个推论是：对每个 $b$，应当存在完全忠实嵌入
\begin{equation*}\label{E: functor A}
\frakA_{J_b}:  \fgRep(J_b(F),\La)\to  \Coh(\Loc_{\CG}\otimes\La),
\end{equation*}
它是把 $\bL_G$ 的拟逆限制到 $i_{b,!}(\fgRep(J_b(F),\La))$ 而得到的。这样的函子之存在，与族中局部朗兰兹的想法密切相关；在 $J_b=G$ 分裂且 $\La$ 是特征零的域的情形，它也在 \cite{Hellmann-Iwahori, BCHN} 中被考虑过。

特别地，对每个开紧子群 $K\subset J_b(F)$，应当存在 $\Loc_{\CG}\otimes\La$ 上的凝聚层
\begin{equation}\label{E: sheaf A}
\frakA_{K,\La}:=\frakA_{J_b}(\delta_{K,\La}),
\end{equation}
使得
\begin{equation}\label{E: Hecke alg isom}
\big(R\End_{\Coh(\Loc_{\CG}\otimes\La)}\frakA_{K,\La}\big)^{\op}\cong \big(R\End_{\Rep(G(F),\La)}\delta_{K,\La}\big)^{\op}=:H_{K,\La}.
\end{equation}
代数 $H_{K,\La}$ 有时称为导出 Hecke 代数，因为（当 $\La=\bZ_\ell$ 或 $\bF_\ell$ 时）它未必集中在上同调次数零。$\frakA_{K,\La}$ 在各种情形下的猜想性描述见 \cite[Sect. 4.3-4.5]{Zhu-Conjecture}。

与整体几何朗兰兹猜想一样，猜想 \ref{C: LLC} 中的等价应当满足一族相容性条件。例如，它应当与两侧的抛物诱导相容，也应当与 $\Shv_c(\frakB(G),\La)$ 上的对偶 $\bD^{\coh}$ 以及 $\Coh(\Loc_{\CG}\otimes\La)$ 的（修正的）Grothendieck--Serre 对偶相容。细节见 \cite{Zhu-Conjecture, Hemo-Zhu-Unipotent}。

另一方面，猜想 \ref{C: LLC} 预测完全复形的范畴 $\Perf(\Loc_{\CG}\otimes\La)$ 作用于 $\Shv_c(\frakB(G),\La)$，这类似于 \S\ref{Global GLC} 中提到的谱作用。\cite{Fargues-Scholze-Geometrization} 的主要结果之一，就是在他们的框架中构造出这样的作用。目前人们还不知道 $\Shv_c(\frakB(G),\La)$ 上是否存在这样的谱作用，但有一些令人信服的证据表明猜想 \ref{C: LLC} 仍应成立。

如前面一样，设 $G$ 可延拓为 $\mO$ 上的约化群。于是有闭子叠
\[\Loc_{\CG}^{\ur}\subset\Loc_{\CG}^{\un}\subset \Loc_{\CG},\]
通常分别称为非分歧参数叠与幂幺参数叠，它们分类那些满足 $\rho(I_F)$ 平凡（分别地，$\rho(I_F)$ 幂幺）的 $\rho$。当 $\La=\bQ_\ell$ 时，$\Loc_{\CG}^{\un}\otimes\bQ_\ell$ 是 $\on{Loc}_{\CG}\otimes\bQ_\ell$ 的一个连通分支。

另一方面，有幂幺子范畴 $\Shv_c^{\un}(\kot(G),\bQ_\ell)\subset \Shv_c(\kot(G),\bQ_\ell)$，粗略地说它是所有 $b\in B(G)$ 的 $J_b(F)$ 幂幺表示范畴 $\fgRep^{\un}(J_b(F),\bQ_\ell)$（由 \cite{Lusztig-Unipotent} 引入）的粘合。下述定理取自 \cite{Hemo-Zhu-Unipotent}，由对定理 \ref{T: Bezrukavnikov equivalence} 取 Frobenius 扭的范畴迹得到。

\begin{thm}\label{T: tame local Langlands}
设 $G$ 是 $\mO$ 上的约化群，并固定钉 $(B,T,e)$。则存在典范等价
\[
\bL^{\un}_G:  \Coh(\Loc_{\CG}^{\un}\otimes\bQ_\ell)\cong  \Shv_{c}^{\un}(\frakB(G),\bQ_\ell).
\]
\end{thm}

对算术应用而言，在此等价下把具体对象对应起来是重要的。我们给出几个例子，更多的对应见 \cite{Hemo-Zhu-Unipotent}（并见 \cite[Sect. 4.3-4.5]{Zhu-Conjecture}）。

\begin{ex}
等价 $\bL_G^{\un}$ 对所有抛孔子群 $K\subset G(F)$（Bruhat--Tits 意义下）给出 \eqref{E: sheaf A} 中的猜想凝聚层，使得 \eqref{E: Hecke alg isom} 成立。例如 $\frakA_{G(\mO),\bQ_\ell}\cong \mO_{\Loc_{\CG}^{\ur}}\otimes\bQ_\ell$，由此得到
\begin{equation}\label{E: derived Satake}
\big(R\End_{\Coh(\Loc_{\CG})}\mO_{\Loc_{\CG}^{\ur}}\big)^{\op}\otimes\bQ_\ell\cong \big(R\End\ \delta_{G(\mO),\bQ_\ell}\big)^{\op}= H_{G(\mO),\bQ_\ell}\cong H_{G(\mO)}^{\cl}\otimes\bQ_\ell.
\end{equation}
由于 $\Loc_{\CG}^{\ur}\cong (\CG|_{d=(q,\sigma)})/\hat{G}$，取零次上同调就还原出 Satake 同构 \eqref{Sat isom}。此外，它还蕴含左端没有高阶上同调，而这一点并不显然。需要指出，\cite[Sect. 4.3]{Zhu-Conjecture} 猜测 $\frakA_{G(\mO),\bZ_\ell}\cong \mO_{\Loc_{\CG}^{\ur}}$，因此 \eqref{E: derived Satake} 中第一个同构应当在 $\bZ_\ell$ 上成立，这称为（猜想性的）导出 Satake 同构。（但一般地 $H_{G(\mO),\bZ_\ell}\neq H_{G(\mO)}^{\cl}\otimes\bZ_\ell$。）

对 $\frakA_{I,\bQ_\ell}$ 还有一个纯 Galois 一侧的描述，即 \cite{BCHN, Zhu-Conjecture} 中定义的幂幺凝聚 Springer 层（并见 \cite{Hellmann-Iwahori}）。
\end{ex}

\begin{ex}\label{Ex: trace of Satake}
由构造，存在 $\bZ_\ell$ 上的叠的态射 $\Loc_{\CG}\to \bB\hat{G}$。设 $\hat{G}$ 在有限投射 $\La$-模上有表示，把它看作 $\bB\hat{G}_\La$ 上的向量丛，以 $\widetilde V$ 记其拉回至 $\Loc_{\CG}\otimes\La$；对 $?=\ur$ 或 $\un$，以 $\widetilde{V}^{?}\in \Perf(\Loc_{\CG}^{?}\otimes\La)$ 记它到 $\Loc_{\CG}^?\otimes\La$ 的限制。注意当 $\La=\bQ_\ell$ 时 $\widetilde{V}^{\ur}\cong \widetilde{V}\otimes \frakA_{G(\mO),\bQ_\ell}$。我们有
\[
\bL_G^{\un}(\widetilde{V}^{\ur})\cong \on{Nt}_!r^!\Sat(V)=:\mS_V,
\]
其中 $r$ 与 $\on{Nt}$ 是下述对应中的态射：
\[
\Hk_G=L^+G\backslash LG/L^+G\xleftarrow{r} LG/\Ad_\sigma L^+G\xrightarrow{\on{Nt}} LG/\Ad_\sigma LG=\kot(G).
\]
特别地，对 $\hat{G}$ 的两个表示 $V$ 与 $W$，存在与复合相容的态射
\begin{equation}\label{E: local S-operator}
R\Hom_{\Loc_{\CG}^{\ur}\otimes\bQ_\ell}(\widetilde{V}^{\ur}, \widetilde{W}^{\ur})\to R\Hom_{\Shv_c(\kot(G),\bQ_\ell)}(\mS_V,\mS_W).
\end{equation}
这样的映射最初由 \cite{Xiao-Zhu-Cycle,Zhu-CDM} 构造，其（非导出 $\Hom$ 空间的）版本随后在 \cite{Yu-Jacquet-Langlands-Higher-Weights} 中被推广到 $\bZ_\ell$ 系数。它有重要的算术应用，将在 \S\ref{Application} 中说明。
\end{ex}

\begin{rmk}
定理 \ref{T: tame local Langlands} 多半可以通过对定理 \ref{T: Monodromic Bezrukavnikov equivalence} 中的等价取 Frobenius 扭的范畴迹，推广到驯顺层。另一方面，如注 \ref{R: int Bez} 所述，把这些等价推广到 $\bZ_\ell$ 系数是重要的。
\end{rmk}

% ============================================================
%  ch02b  §2.2 函数域上的整体算术朗兰兹
%  原文对应：ICM-arXiv.tex 937--1009 行
%  编号：定理 2.2.1；注 2.2.2；公式 (2.5)；猜想 2.2.3；
%        定理 2.2.4；推论 2.2.5
% ============================================================

\subsection{函数域上的整体算术朗兰兹}\label{S: Global Langlands function fields}

接下来转向算术朗兰兹对应的整体方面。如开头所述，其经典表述粗略地预测（不可约）Galois 表示的集合与（尖点）自守表示的集合之间存在自然的对应。与局部情形一样，除 $\GL_n$ 情形（这是 \cite{Lafforgue-GLn} 的定理）之外，这样的表述并不容易精确化。另一方面，\S\ref{Global GLC} 中的整体几何朗兰兹猜想以及去范畴化／迹的哲学提示：整体算术朗兰兹可以、而且多半应当表述为两个向量空间之间的同构，这两个空间分别来自 Galois 一侧与自守一侧。本小节在整体函数域情形表述这样一个猜想。

设 $F=\bF_q(X)$ 是 $\bF_q$ 上几何连通光滑射影曲线 $X$ 的函数域。以 $\eta=\Spec F$ 记 $X$ 的生成点，$\overline\eta$ 记 $\eta$ 上的一个几何点。以 $|X|$ 记 $X$ 的闭点集。对 $v\in |X|$，以 $\mO_v$ 记 $X$ 在 $v$ 处的完全局部环，$F_v$ 为其分式域。令 $\bO_F=\prod_{v\in|X|}\mO_v$ 为整 adèle 环，$\bA_F=\prod'_{v\in|X|}F_v$ 为 adèle 环。对有限非空的位置集 $Q$，以 $W_{F,Q}$ 记 $F$ 的在 $Q$ 外非分歧的 Weil 群。

设 $G$ 是 $F$ 上的连通约化群。与局部情形类似，表述整体猜想的第一步是 \cite{Zhu-Conjecture} 中的下述定理。

\begin{thm}
设 $\ell\nmid 2p$。把 $\bZ_\ell$-代数 $A$ 送到满足 $d\circ\rho=(\cycl^{-1},\pr)$ 的（强）连续同态 $\rho: W_{F,Q}\to \CG(A)$ 所构成空间的预叠，可由一个导出概形 $\Loc_{\CG,Q}^\Box$ 表示；该概形是若干导出仿射概形的无交并，每个在 $\bZ_\ell$ 上平坦、有限型。若 $Q\neq\emptyset$，则 $\Loc_{\CG,Q}^\Box$ 拟光滑。
\end{thm}

于是定义整体朗兰兹参数叠为 $\Loc_{\CG,Q}=\Loc_{\CG,Q}^\Box/\hat{G}$。与局部情形类似（见例 \ref{Ex: trace of Satake}），对 $\hat{G}_\La$ 在有限投射 $\La$-模上的表示，把它看作 $\bB\hat{G}_\La$ 上的向量丛，以 $\widetilde V$ 记其拉回至 $\Loc_{\CG,Q}\otimes\La$。若 $V$ 是 $(\CG)^S$ 的某个表示沿对角嵌入 $\hat{G}\to (\CG)^S$ 的限制，则 $\widetilde V$ 上有自然的（强）连续 $W_{F,Q}^S$-作用（见 \cite[Sect. 2.4]{Zhu-Conjecture}）。对 $F$ 的位置 $v$，以 $\Loc_{\CG,v}$ 记 $G_{F_v}$ 的局部朗兰兹参数叠。令
\[
\res: \Loc_{\CG,Q}\to \prod_{v\in Q}\Loc_{\CG,v}
\]
为由整体参数限制到局部参数所得的映射（它由 $W_{F_v}\to W_{F,Q}$ 诱导）。下文将考虑沿此映射对 $\prod_{v\in Q}\Loc_{\CG,v}$ 上凝聚层取 $!$-拉回。

\begin{rmk}
(1) 事实上，当 $Q=\emptyset$ 时 $\Loc_{\CG,Q}$ 的定义需要略作修改。

(2) 与局部情形不同，$\Loc_{\CG,Q}$ 一般带有非平凡的导出结构（见 \cite[Rem. 3.4.5]{Zhu-Conjecture}）。以 ${}^{cl}\Loc_{\CG,Q}$ 记其底层的经典叠。

(3) \cite{AGKRRV} 给出了 $\Loc_{\CG,Q}\otimes\bQ_\ell$ 的另一种定义。
\end{rmk}

接着转到自守一侧。为简单起见，本小节设 $G$ 在 $\bF_q$ 上分裂。固定一个水平，即开紧子群 $K\subset G(\bO_F)$。设 $Q$ 是由那些满足 $K_v\neq G(\mO_v)$ 的 $v$ 构成的位置集。对有限集 $S$，以 $\Sht_K(G)_{(X-Q)^S}$ 记 $(X-Q)^S$ 上由 $X$ 上带 $K$-水平结构的 $S$-leg（legs）落在 $X-Q$ 中的 $G$-shtuka 的模空间所给出的 ind-Deligne--Mumford 叠。（该模空间的基本构造与性质例如见 \cite{Lafforgue-Automorphic-to-Galois}。）它沿对角映射 $\overline\eta\to (X-Q)\xrightarrow{\Delta}(X-Q)^S$ 的基变换记作 $\Sht_K(G)_{\Delta(\overline\eta)}$。

对 $(\CG)^S$ 在有限投射 $\La$-模上的每个表示 $V$，几何 Satake \eqref{factorizable Sat}（把 $D$ 换成 $X-Q$，并允许 $\La=\bZ_\ell$）给出 $\Sht_K(G)_{(X-Q)^S}$ 上的反常层 $\on{Sat}_S(V)$。以 $C_c\bigl(\Sht_K(G)_{\Delta(\overline\eta)},\Sat_S(V)\bigr)$ 记 $\Sht_K(G)_{\Delta(\overline\eta)}$ 以 $\Sat_S(V)$ 为系数的总的紧支撑上同调（的上链复形）。它带有（强）连续的 $W_{F,Q}^S$-作用（导出层次上这种作用的构造见 \cite{HRS-Kunneth}，其基础是 \cite{Xue-Finiteness, Xue-Smoothness}），还带有相应的整体（导出）Hecke 代数（系数在 $\La$ 中）
\begin{equation}\label{E: derived hecke global}
H_{K,\La}=\Big(R\End \big(c\on{-ind}_K^{G(\bA_F)}(\La)\big)\Big)^{\op}
\end{equation}
的作用。例如，若 $V=\mathbf{1}$ 是平凡表示，则（在 $G$ 分裂的假设下）
\[C_c\bigl(\Sht_K(G)_{\Delta(\overline\eta)}, \on{Sat}_{\{1\}}(\mathbf{1})\bigr)= C_c(G(F)\backslash G(\bA)/K,\La).\]
这里 $G(F)\backslash G(\bA)/K$ 看作 $\overline\eta$ 上的离散 DM 叠，而 $C_c(G(F)\backslash G(\bA)/K,\La)$ 记其紧支撑上同调。当 $\La=\bQ_\ell$ 时，这就是 $G(F)\backslash G(\bA)/K$ 上紧支撑函数构成的空间。

我们将固定 $G$ 的一个钉 $(B,T,e)$ 以及非退化特征 $\psi: F\backslash \bA\to \bZ_\ell[\mu_{p}]^\times$，它按猜想 \ref{C: LLC} 对每个 $v\in Q$ 给出猜想性的等价 $\bL_{v}$。特别地，与 $K_v\subset G(F_v)$ 相对应，有 $\Loc_{\CG,v}$ 上的猜想凝聚层 $\frakA_{K_v}$（见 \eqref{E: sheaf A}）。

\begin{conj}\label{C: gen coh}
存在自然的 $(W_{F,Q}^S\times H_{K,\La})$-等变同构
\[
R\Gamma\bigl(\on{Loc}_{\CG,Q}\otimes\La, \widetilde{V}\otimes  \res^!(\boxtimes_{v\in Q} \frakA_{K_v})\bigr)\cong C_c\bigl(\on{Sht}_K(G)_{\Delta(\overline\eta)}, \on{Sat}_S(V)\bigr).
\]
\end{conj}

该猜想更一般的形式（允许“广义水平结构”）以及在各种特殊情形下的例子，见 \cite[Sect. 4.7]{Zhu-Conjecture}。这个猜想可以看作函数域上整体朗兰兹对应的精确形式：它给出了把 Galois 表示与自守表示对应起来的精确配方。（例如，V. Lafforgue 的巡游算子就编码在此同构之中，见下文。）此外，这样的同构符合 Arthur--Kottwitz 重数公式，同时把该公式推广到整系数层次，从而与自守提升理论相关联。

该猜想最令人信服的证据是下述定理 \cite{Lafforgue-Zhu-Elliptic, Zhu-Conjecture}，它正是 Drinfeld 对 Lafforgue 构造所作解释（在启发式层次上）所提示的。

\begin{thm}\label{T: gen LZ}
对每个 $i$，存在 ${}^{cl}\Loc_{\CG,Q}\otimes\bQ_\ell$ 上的拟凝聚层 $\frakA^i_K$，带有 $H_{K,\bQ_\ell}$ 的作用，使得对 $(\CG)^S$ 的每个有限维 $\bQ_\ell$-表示 $V$，都有自然的 $(W_{F,Q}^S\times H_{K,\bQ_\ell})$-等变同构
\begin{equation*}\label{E: coh Sht}
\Gamma\big({}^{cl}\Loc_{\CG,Q}\otimes\bQ_\ell,\widetilde V\otimes \frakA^i_K\big)\cong H^i_c\bigl(\Sht_K(G)_{\Delta(\overline\eta)},\Sat_S(V)\bigr).
\end{equation*}
\end{thm}

需要指出，\cite{Lafforgue-Zhu-Elliptic, Zhu-Conjecture} 实际上对任意 $G$ 都证明了这个定理。此外，当 $K$ 处处超特殊时，\cite{AGKRRV} 证明 \eqref{E: coh Sht} 在导出层次上成立。

同构 \eqref{E: coh Sht} 诱导出 $\Gamma\big({}^{cl}\Loc_{\CG,Q}\otimes\bQ_\ell,\mO\big)$ 在右端的作用。这正是 V. Lafforgue 的巡游算子所给出的作用，它诱导出右端（特别是 $C_c(G(F)\backslash G(\bA_F)/K, \bQ_\ell)$）按半单朗兰兹参数的分解。如 \cite{Lafforgue-Zhu-Elliptic} 所说明，在椭圆朗兰兹参数上，这一同构与 Arthur--Kottwitz 重数公式密切相关。当 $G=\GL_n$ 时，它给出下述推论，推广了 \cite{Lafforgue-GLn}。

\begin{cor}
设 $\pi$ 是 $\GL_n$ 的尖点自守表示，其相伴的不可约 Galois 表示 $\rho_\pi: W_{F,Q}\to \GL_n(\La)$ 取值于某个有限扩张 $\La/\bQ_\ell$，并以 $\frakm_\pi$ 记 $\Gamma\big({}^{cl}\Loc_{\CG,Q}\otimes\La,\mO\big)$ 中相应的极大理想。则存在 $(W_{F,Q}^S\times H_K)$-等变同构
\[
H^*_c\big(\Sht_K(G)_{\Delta(\overline\eta)}, \Sat_S(V)\big)/\frakm_\pi\cong  V_{\rho_\pi}\otimes \pi^K.
\]
特别地，左端只集中在上同调次数零。
\end{cor}

% ============================================================
%  ch02c  §2.3 Jacquet--Langlands 提升的几何实现
%  原文对应：ICM-arXiv.tex 1011--1069 行
%  编号：定义 2.3.1；猜想 2.3.2；公式 (2.6)(2.7)(2.8)；定理 2.3.3
% ============================================================

\subsection{Jacquet--Langlands 提升的几何实现}\label{geometric JL}

数域上的整体朗兰兹对应要复杂得多。事实上，理论中有分析的部分，目前看来无法纳入范畴化／去范畴化的框架。即便只限于理论的代数／算术部分，也会有来自 $\ell$ 处与 $\infty$ 处的复杂性。特别是，$\ell$ 处与 $\infty$ 处的局部朗兰兹对应的范畴形式尚未被完全理解。

尽管如此，在与 Emerton 以及 Emerton--Gee 的合作文章（\cite{Emerton-Zhu, Emerton-Gee-Zhu}）中，我们将对志村簇的上同调、乃至一般局部对称空间的上同调，提出猜想性的 Galois 理论描述，与猜想 \ref{C: gen coh} 平行。本小节只回顾 \cite{Zhu-Conjecture} 中关于借助志村簇上同调几何地实现 Jacquet--Langlands 提升的一个猜想，并讨论 \cite{Xiao-Zhu-Cycle, Hemo-Zhu-Unipotent} 中走向该猜想的若干结果。

先固定一些记号与假设，本小节固定一个素数 $p$。以 $\bA_f=\prod'_q\bQ_q$ 记 $\bQ$ 的有限 adèle 环，$\bA_f^p=\prod'_{q\neq p}\bQ_p$。记 $\overline\eta=\Spec \overline{\bQ}$，其中 $\overline{\bQ}$ 是 $\bQ$ 在 $\bC$ 中的代数闭包。对志村数据 $(G,X)$，以 $\mu$ 记由 $(G,X)$ 按通常方式确定的 $\hat{G}$（关于 $(\hat{B},\hat{T})$）的（微小的）支配权，以 $V_\mu$ 记 $\hat{G}$ 的以 $\mu$ 为最高权的微小表示。以 $E\subset\overline{\bQ}\subset\bC$ 记 $(G,X)$ 的反射域，并记 $d_\mu=\dim X$。

对水平（即开紧子群）$K=K_pK^p\subset G(\bQ_p)G(\bA_f^p)$，以 $\Sh_K(G)$ 记相应水平的志村簇（定义在反射域 $E$ 上），以 $\Sh_K(G)_{\overline\eta}$ 记它沿 $E\to \overline\bQ$ 的基变换。设 $v$ 是 $E$ 上位于 $p$ 之上的位置。所谓特化 $\on{sp}: \overline\eta\to \overline{v}$，指从 $\overline\eta$ 到 $\mO_E$ 在 $v$ 处的严格 Hensel 化的态射。

为避免 Galois 上同调带来的诸多复杂情形（例如扩张纯内形式与内形式之别），也避免几何上的一些复杂情形（例如志村簇与 Shtuka 模空间之间的关系），本小节余下部分设 $G$ 为伴随型；一般 $G$ 的情形见 \cite{Xiao-Zhu-Cycle}，对 $G$ 限制更少的版本见 \cite{Zhu-Conjecture}。

\begin{dfn}\label{D: prime-to-p trivial}
设 $G$ 是 $\bQ$ 上的连通约化群。$G$ 的一个 $p$-外（分别地，有限）平凡化内形式，指 $\bQ$ 上的一个 $G$-挠子 $\beta$，并带有 $\beta$ 在 $\bA_f^p$ 上（分别地，在 $\bA_f$ 上）的平凡化。此时 $G':=\Aut(\beta)$ 是 $G$ 的内形式（因此 $G$ 与 $G'$ 的对偶群被典范地等同），并带有同构 $\theta: G(\bA^p_f)\cong G'(\bA^p_f)$（分别地，$\theta: G(\bA_f)\cong G'(\bA_f)$）。
\end{dfn}

今设 $(G,X)$ 与 $(G',X')$ 是两个志村数据，$G'$ 是 $G$ 的 $p$-外平凡化内形式。借助 $\theta$，可把 $K^p\subset G(\bA_f^p)$ 搬为开紧子群 ${K'}^p\subset G'(\bA_f^p)$。我们把 $p$-外（导出）Hecke 代数 $H_{K^p,\La}$（定义方式与 \eqref{E: derived hecke global} 相同）与 $H_{{K'}^p,\La}$ 等同，并统记作 $H_{K^p,\La}$。设 $K'_p\subset G'(\bQ_p)$ 是开紧子群，记 $K'=K'_p{K'}^p$ 为相应水平。

固定 $G_{\bQ_p}$ 与 $G'_{\bQ_p}$ 的一个拟分裂内形式 $G^*_{\bQ_p}$，并带有钉 $(B^*_{\bQ_p},T^*_{\bQ_p},e^*)$；把 $G_{\bQ_p}$ 实现为 $J_b$、把 $G'_{\bQ_p}$ 实现为 $J_{b'}$，其中 $b,b'\in B(G^*_{\bQ_p})$。在 $G$ 与 $G'$ 为伴随型的假设下，这样的 $b,b'$ 存在且唯一。于是按 \eqref{E: sheaf A}，在 $G^*_{\bQ_p}$ 于 $\La$ 上的局部朗兰兹参数叠 $\Loc_{\CG,p}\otimes\La$ 上有猜想凝聚层 $\frakA_{K_p,\La}$ 与 $\frakA_{K'_p,\La}$。

\begin{conj}\label{C: geo JL}
对特化映射 $\on{sp}: \overline{\eta}\to \overline{v}$ 的任一选取，存在自然映射
\begin{multline}\label{E: geoJL}
R\Hom_{\Coh(\Loc_{\CG,p}\otimes\La)}\bigl(\widetilde{V_\mu}\otimes \frakA_{K_p,\La}, \widetilde{V_{\mu'}}\otimes \frakA_{K'_p,\La}\bigr)\\
\to R\Hom_{H_{K^p,\La}}\bigl(C_c(\Sh_{K}(G)_{\overline\eta}, \La[d_{\mu}]), C_c(\Sh_{K'}(G')_{\overline\eta},\La[d_{\mu'}])\bigr),
\end{multline}
并与复合相容。特别地，有（$E_1$-）代数同态
\begin{equation}\label{E: S operator}
S:R\End_{\Coh(\Loc_{\CG,p}\otimes\La)}\bigl( \widetilde{V_\mu}\otimes \frakA_{K_p,\La}\bigr)\to R\End_{H_{K^p,\La}}\bigl(C_c(\Sh_{K}(G)_{\overline\eta},\La)\bigr),
\end{equation}
它与 \eqref{E: geoJL} 相容。此外，诱导出的作用
\begin{equation}\label{E: S action}
H_{K_p,\La}\stackrel{\eqref{E: Hecke alg isom}}{\cong} R\End(\frakA_{K_p,\La})\to R\End\bigl(\widetilde{V_\mu}\otimes \frakA_{K_p,\La}\bigr)\xrightarrow{S} R\End_{H_{K^p,\La}}\bigl(C_c(\Sh_{K}(G)_{\overline\eta},\La)\bigr)
\end{equation}
与 $H_{K_p,\La}$ 在 $C_c(\Sh_{K}(G)_{\overline\eta},\La)$ 上的自然 Hecke 作用一致（因此与特化映射 $\on{sp}$ 无关）。
\end{conj}

如果对 $C_c(\Sh_{K}(G)_{\overline\eta},\La)$ 有类似于猜想 \ref{C: gen coh} 的 Galois 理论描述，本猜想便是其推论；但本猜想的表述并不需要 $\bQ$ 的整体朗兰兹参数叠存在。无论如何，走向 $C_c(\Sh_{K}(G)_{\overline\eta},\La)$ 的 Galois 理论描述，或许需要以猜想 \ref{C: geo JL} 为输入。还要指出，与函数域情形一样，\cite[Sect. 4.7]{Zhu-Conjecture} 中有本猜想的更一般版本，允许“广义水平结构”，从而使 Igusa 簇的上同调也能出现。

下述定理在特殊情形下验证了该猜想。

\begin{thm}
\label{introT: spectral action}
设志村数据 $(G,X)$ 与 $(G',X')$ 都是阿贝尔型，且 $G'$ 是 $G$ 的有限平凡化内形式。设 $G_{\bQ_p}$ 非分歧（从而 $G'_{\bQ_p}$ 也非分歧）。
\begin{enumerate}
\item\label{introT: spectral action-1} 当 $\La=\bQ_\ell$ 且 $K_p\subset G(\bQ_p)$、$K'_p\subset G'(\bQ_p)$ 为抛孔子群（Bruhat--Tits 意义下）时，映射 \eqref{E: geoJL}（从而 \eqref{E: S operator}）存在。
\item\label{introT: spectral action-2} 若 $K_p$ 超特殊，则当 $\La=\bZ_\ell$ 时映射 \eqref{E: geoJL}（从而 \eqref{E: S operator}）存在，至少对非导出的 $\Hom$ 空间如此。此外，$H_{K_p}^{\cl}$ 通过 \eqref{E: S action} 在 $H_c^*(\Sh_{K}(G)_{\overline\eta},\La)$ 上的作用与 $H_{K_p}^{\cl}$ 的自然作用一致。
\end{enumerate}
\end{thm}

第 \eqref{introT: spectral action-1} 部分在 \cite{Xiao-Zhu-Cycle, Hemo-Zhu-Unipotent} 中证明。证明包含两个要素：一是 \cite{Xiao-Zhu-Cycle} 对 $\Sh_K(G)$ 与 $\Sh_{K'}(G')$ 的 mod $p$ 纤维之间物理对应的构造（这里用到了目前需要假设 $G$ 与 $G'$ 在 $p$ 处非分歧这一点）；二是定理 \ref{T: tame local Langlands}（因此需要 $\La=\bQ_\ell$）。当 $K_p$ 超特殊时，可以取 $\bZ_\ell$ 系数，因为 \eqref{E: local S-operator} 的（非导出版本）对 $\bZ_\ell$ 系数成立，这得益于 \cite{Yu-Jacquet-Langlands-Higher-Weights}。事实上，此时还可以允许志村簇上有非平凡局部系统（见 \cite{Yu-Jacquet-Langlands-Higher-Weights}）。

最后一个断言即志村簇情形的 $S=T$。当 $d_{\mu}=\dim\Sh_K(G)=0$ 时它包含于 \cite{Xiao-Zhu-Cycle}；一般情形在 \cite{Wu-S=T, Zhu-S=T} 中借助 \cite{Scholze-Berkeley, Fargues-Scholze-Geometrization} 的基础性工作得到证明。

% ============================================================
%  ch03a  §3 引言 + §3.1 志村簇的局部模型
%  原文对应：ICM-arXiv.tex 1070--1105 行
%  编号：公式 (3.1)；定理 3.1.1；注 3.1.2
% ============================================================

\section{对算术几何的应用}\label{Application}

除前面提到的几何朗兰兹（的想法）对经典朗兰兹纲领的直接应用之外，我们再讨论一些算术应用，它们大多与志村簇有关，且基于作者本人的工作。需要指出，几何朗兰兹（的想法）对算术问题还有许多其他值得注意的应用，例如 \cite{Gaitsgory-deJong, HNY-Kloosterman, Yun-motive-exceptional-group, Xu-Zhu-Bessel, LLLM}，这里仅举几例。

\subsection{志村簇的局部模型}\label{S: local model}

志村簇整模型（带抛孔子群水平）的理论，始于（隐含在 Kronecker 工作中的）对带 $\Ga_0(p)$-水平的椭圆模曲线的 mod $p$ 约化的理解。我们讨论这一理论的一小部分，即借助局部模型理论研究这些整模型的 étale 局部结构。局部模型理论的近期发展深受几何朗兰兹纲领的影响。

沿用 \S\ref{geometric JL} 中关于志村簇的记号（但本小节不假设 $G$ 为伴随型）。设 $(G,X)$ 为志村数据，$K$ 为所选水平，满足 $K_p=\mG(\bZ_p)$，其中 $\mG$ 是 $G_{\bQ_p}$ 在 $\bZ_p$ 上的某个抛孔子群概形（Bruhat--Tits 意义下）。于是对 $E$ 上位于 $p$ 之上的位置 $v$，局部模型图指 $\mO_{E_v}$ 上拟射影概形之间的对应
\begin{equation}\label{E: local model diagram}
\scrS_K\leftarrow\widetilde{\scrS}_K\xrightarrow{\tilde\varphi} M^{\loc}_{\mG},
\end{equation}
其中 $\scrS_K$ 是 $\Sh_K(G)$ 在 $\mO_{E_v}$ 上的整模型，$\widetilde{\scrS}_K$ 是 $\scrS_K$ 上的 $\mG_{\mO_{E_v}}$-挠子，$M^{\loc}_{\mG}$ 是所谓的局部模型，即 $\mO_{E_v}$ 上带 $\mG_{\mO_{E_v}}$-作用的平坦射影概形，而 $\tilde{\varphi}$ 是相对维数为 $\dim G$ 的 $\mG_{\mO_{E_v}}$-等变光滑态射。因此 $M_{\mG}^{\loc}$ 刻画了 $\scrS_K$ 的 étale 局部结构。另一方面，$M_{\mG}^\loc$ 上 $\mG_{\mO_{E_v}}$-作用的存在，使它比 $\scrS_K$ 更易研究。

局部模型最初的构造，基于把抛孔子群概形实现为（$F$ 上除代数中）带双线性型的向量空间里一条自对偶格链的稳定化群（的幺单位连通分支），综述与文献见 \cite{PRS-Survey}。这一途径多少有些特设，且局限于所谓的（P）EL（局部）志村数据。基于 \S\ref{Tamely local} 中叠 $\Hk_{\mG,D}$ 的 $\bZ_p$-类比之构造的新途径，由 \cite{Pappas-Zhu-Local-Model} 系统引入（其中假设 $G$ 驯顺，该假设后来在 \cite{Levin-Weil-Restriction, Lourenco} 中被去掉）。在该文中，这种 $\bZ_p$-类比（更确切地说，是相应的 Beilinson--Drinfeld 型仿射 Grassmann 簇 $\Gr_{\mG,\bZ_p}$）的构造，基于对某个“二维抛孔子”群概形 $\widetilde\mG$ 的构造：它是 $\bZ_p[\varpi]$ 上的群概形，沿 $\bZ_p[\varpi]\xrightarrow{\varpi\mapsto p}\bZ_p$ 的限制可还原出 $\mG$。（综述见 \cite{Zhu-ICCM}。）

这种仿射 Grassmann 簇的另一 $p$-进版本 $\Gr_{\mG,\Spd\bZ_p}$，在 \cite{Scholze-Berkeley} 中于解析的完满世界中给出了更直接的构造。无论哪种情形，局部模型都定义为生成纤维中对应于 $\mu$ 的 Schubert 簇的平坦闭包。此外，近期工作 \cite{AGLR} 表明这两种构造一致。

下述定理取自 \cite{AGLR}，是关于局部模型存在性及其性质的最新结果。

\begin{thm}
设 $G$ 是 $p$-进域 $F$ 上的连通约化群。除 $p=2$ 的奇酉情形与 $p=3$ 的三对偶情形之外，对 $G$ 在 $\mO$ 上的任一抛孔子群概形 $\mG$，以及 $G$ 的、定义在 $F$ 的有限扩张 $E/F$ 上的微小余权共轭类 $\mu$，都存在 $\mO_E$ 上的正规平坦射影概形 $M_{\mG,\mu}^{\loc}$，带有 $\mG_{\mO_E}$-作用，使得 $M_{\mG,\mu}^{\loc}\otimes E$ 与 $G_E$ 中对应于 $\mu$ 的部分旗簇 $\Fl_{\mu}$ 是 $G_E$-等变同构的，而 $M_{\mG}^{\loc}\otimes k_E$ 与 $LG/L^+\mG\otimes k_E$ 中 $\mu$-容许集上的 Schubert 簇之并（的典范去完美化）是 $(\mG\otimes k_E)$-等变同构的。此外，$M_{\mG}^{\loc}$ 正规、Cohen--Macaulay，且其特殊纤维的每个几何不可约分支都正规且 Cohen--Macaulay。
\end{thm}

本小节以几条注记结束。

\begin{rmk}
\begin{enumerate}
\item 一旦建立局部模型图 \eqref{E: local model diagram}，该定理也就给出志村簇整模型的相应性质。

\item 研究局部模型特殊纤维的一个关键要素，是 Pappas--Rapoport 的凝聚猜想 \cite{Pappas-Rapoport-Ramified-Unitary}，它由 \cite{Zhu-Coherence-Conjecture} 证明（证明用到了融合积的想法）。

\item 局部模型理论的一个重要动机／应用是 Haines--Kottwitz 猜想 \cite{Haines-Test-Function}：它预测，抛孔 Hecke 代数 $H_{K_p}^{\cl}$ 中的某个中心元应当作为计算 $\Sh_{K}(G)$ 的 Hasse--Weil zeta 函数的迹公式中的测试函数。如 \S\ref{Tamely local} 所述，这一猜想正是 Gaitsgory 中心层构造 \eqref{E: central sheaf} 的动机之一。在构造出 $\bZ_p$ 上的局部 Hecke 叠 $\Hk_{\mG,\bZ_p}$（无论用 \cite{Pappas-Zhu-Local-Model} 还是 \cite{Scholze-Berkeley} 的版本）之后，就可以在混合特征下模仿构造 \eqref{E: central sheaf}，从而解决 Kottwitz 猜想。最新结果同样见 \cite{AGLR}。
\end{enumerate}
\end{rmk}

% ============================================================
%  ch03b  §3.2 同余关系
%  原文对应：ICM-arXiv.tex 1107--1140 行
%  编号：公式 (3.2)；定理 3.2.1；定理 3.2.2
% ============================================================

\subsection{同余关系}\label{S: Congruence}

沿用 \S\ref{geometric JL} 中关于志村簇的记号，并（为简单起见）保留其中的假设。设 $(G,X)$ 是阿贝尔型志村数据，$K$ 是满足 $K_p$ 超特殊的水平。设 $v\mid p$ 是 $E$ 的位置。则 $\Sh_K(G)$ 有定义在 $\mO_{E,(v)}$ 上的典范整模型 $\scrS_K$（\cite{Kisin-Integral-Model}）。

以 $\overline{\scrS}_K$ 记其 mod $p$ 纤维，它是定义在 $v$ 的剩余域 $k_v$ 上的光滑簇。以 $\sigma_v$ 记 $\Ga_{k_v}$ 中的几何 Frobenius。定理 \ref{introT: spectral action} 给出 $\End_{\Loc_{\CG,p}^{\ur}}(\widetilde{V_\mu})$ 在 $H^*_c(\overline\scrS_{K,\overline{k_v}},\bZ_\ell)$ 上的作用，下文将看到它有重要的推论。

同余关系猜想（又称 Blasius--Rogawski 猜想）推广了模曲线的经典 Eichler--Shimura 同余关系 $\Fr_p=T_p+V_p$；它预测：在 $\overline{\scrS}_K\times \overline{\scrS}_K$ 的 Chow 群中，$\overline{\scrS}_K$ 的 Frobenius 自同态满足一个多项式方程，其系数是某些 Hecke 对应的 mod $p$ 约化。定理 \ref{introT: spectral action} 与 \cite[Sect. 6.3]{Xiao-Zhu-Invariants} 一起，在上同调层次上蕴含了这一猜想。

对 ${}^c(G_{\bQ_p})$ 的每个表示 $V$，其特征 $\chi_V$（看作 $\CG|_{d=(p,\sigma_p)}$ 上的 $\hat{G}$-不变函数）通过 Satake 同构 \eqref{Sat isom} 给出元素 $h_V\in H_{G(\bZ_p)}^{\cl}$。

\begin{thm}
下述恒等式
\begin{equation}\label{E: congruence relation}
\sum_{i=0}^{n}(-1)^j h_{\chi_{\wedge ^j V}} \sigma_v^{\dim V-j}=0
\end{equation}
在 $\End (H^*_c(\overline{\scrS}_{K ,\overline{k_v}},\bZ_\ell))$ 中成立，其中 $V=\on{Ind}_{{}^c(G_{E_v})}^{{}^c(G_{\bQ_p})}V_\mu$ 是 $V_\mu$ 的张量诱导\zy{原文 \eqref{E: congruence relation} 中求和指标写作 $i$，而符号与指数处写作 $j$，且未说明上界 $n$ 的取值。此处照录原文，未作改动；疑为 $i$、$j$ 混用之误，已登记于 errata 表。}。
\end{thm}

事实上，由 \cite[Sect. 6.3]{Xiao-Zhu-Invariants}，把 $h_{\chi_{\wedge ^i V}}$ 换成 $S(\chi_{\wedge ^i V})$（其中 $S$ 来自定理 \ref{introT: spectral action} 的 \eqref{introT: spectral action-1}）后等式成立；再由该定理的第 \eqref{introT: spectral action-2} 部分，便可把 $S(\chi_{\wedge ^i V})$ 换回 $h_{\chi_{\wedge ^i V}}$。这条通向 \eqref{E: congruence relation} 的途径，是 V. Lafforgue 对 $\Sht_K(G)$ 的 Eichler--Shimura 关系之处理方式 \cite{Lafforgue-Automorphic-to-Galois} 在志村簇上的类比。传统上还有另一条途径，即直接研究 Hecke 算子的 mod $p$ 约化，Siegel 模簇情形自 \cite{Faltings-Chai-Abelian} 开始；最新进展及相关文献见 \cite{Lee-ES-Relation}。那条途径会在代数对应的层次上给出 \eqref{E: congruence relation}。

今设 $(G,X)=(\Res_{F^+/\bQ}(G_0)_{F^+}, \prod_{\varphi :F^+\to \bR} X_0)$，其中 $(G_0,X_0)$ 是志村数据，$F^+$ 是全实域。如前，取素数 $p$ 使 $K_p$ 超特殊，特别地 $p$ 在 $F^+$ 中非分歧。此外为简单起见，设 $G_{0,\bQ_p}$ 分裂（于是对 $E$ 上位于 $p$ 之上的位置 $v$，有 $E_v=\bQ_p$）。以 $\bF$ 记 $\bF_p$ 的代数闭包。

设 $\{w_i\}_i$ 是 $F^+$ 上位于 $p$ 之上的素数集合，$k_i$ 为 $w_i$ 的剩余域。对每个 $i$，再固定嵌入 $\rho_i: k_i\to\bF$。于是有自然映射
\[
\prod_i(\bZ^{f_i}\rtimes \frakS_{f_i})\to \End_{\Loc_{\CG,p}^{\ur}}(\widetilde{V_\mu}),
\]
其中 $\frakS_{f_i}$ 是 $f_i$ 个文字上的置换群。结合定理 \ref{introT: spectral action}，便得到下述结果（\cite{Xiao-Zhu-Cycle}）。

\begin{thm}
存在 $\prod_i(\bZ^{f_i}\rtimes \frakS_{f_i})$ 在 $H_c^*(\overline\scrS_{K,\overline{\bF}},\bZ_\ell)$ 上的作用，使得 $\sigma_p$ 的作用分解为 $\sigma_p=\prod_i \sigma_{p,i}$，其中 $\sigma_{p,i}=((1,0,\ldots,0), (12\cdots f_i))\in \bZ^{f_i}\rtimes\frakS_{f_i}$。每个 $\sigma_{p,i}^{f_i}$ 都满足一个类似于 \eqref{E: congruence relation} 的多项式方程。
\end{thm}

该定理在某种程度上反映了 Nekovář--Scholl 的 plectic 上同调猜想 \cite{Nekovar-Scholl-Plectic}。

% ============================================================
%  ch03c  §3.3 志村簇 mod p 纤维上的一般 Tate 圈
%  原文对应：ICM-arXiv.tex 1144--1245 行
%  编号：定理 3.3.1；注 3.3.2；例 3.3.3；公式 (3.3)(3.4)
% ============================================================

\subsection{志村簇 mod $p$ 纤维上的一般 Tate 圈}\label{S: Tate}

在 \cite{Xiao-Zhu-Cycle} 中，我们应用定理 \ref{introT: spectral action} 验证了许多志村簇的 mod $p$ 纤维上 Tate 猜想的“一般”情形。沿用 \S\ref{S: Congruence} 的记号，并（为简单起见）保留其中的假设。以 $(\overline{\mathscr S}_{K,\overline{k_v}})^\pf$ 记 $\overline{\mathscr S}_{K,\overline{k_v}}$ 的完美化（即把它看作 $\Aff_{\overline{k_v}}^{\pf}$ 上的完美预层）；对 $\overline{\mathscr S}_{K,\bar{k}}$ 的每一点附上一个带 $G$-结构的 $F$-同晶体（见 \cite{Kisin-Integral-Model, Xiao-Zhu-Cycle}），便可定义所谓的 Newton 映射
\[
\on{Nt}: (\overline{\mathscr S}_{K,\overline{k_v}})^\pf\to \kot(G_{\bQ_p})_{\overline{k_v}}.
\]
于是 $\kot(G_{\bQ_p})_{\overline{k_v}}$ 的 Newton 分层（见 \S\ref{S: Local Langlands}）诱导出 $\overline{\mathscr S}_{K,\overline{k_v}}$ 被局部闭子簇分层。已知 $\on{Nt}$ 的像含唯一基本元 $b$，$\overline{\mathscr S}_{K,\overline{k_v}}$ 中相应的子簇是闭的，称为基本 Newton 层，记作 $\overline{\scrS}_b$。

设 $m$ 是几何 Frobenius $\sigma_p$ 在 $\xch(\hat{T})$ 上作用的阶。令
\[
\Lambda^\Tate_p=\Big\{\la\in \xch(\hat T)\; \Big|\; \sum_{i=0}^{m-1} \sigma_p^i(\la)=0\Big\}\subset \xch(\hat T).
\]
对 $\hat G_{\bQ_\ell}$ 的表示 $V$ 与 $\la\in\xch(\hat T)$，以 $V(\la)$ 记 $V$ 的 $\la$-权子空间（关于 $\hat T$），并令
\[
V^\Tate=\bigoplus_{\la\in\Lambda^\Tate_p} V(\la).
\]
我们特别关心条件 $V^\Tate_\mu\neq 0$。如 \cite{Xiao-Zhu-Cycle} 的引言所说明，在志村簇上同调的猜想性 Galois 理论描述（类似于猜想 \ref{C: gen coh}）之下，对一个在 $p$ 处的 Satake 参数足够一般的 Hecke 模 $\pi_f$，该向量空间维数的某个倍数 $a(\pi_f)$ 应等于 $\overline{\scrS}_{K,\overline{k_v}}$ 中间维数紧支撑上同调的 $\pi_f$-分支中 Tate 类空间的维数。此外，这个空间通常很大。例如，当 $G$ 是二次虚域上符号为 $(i,n-i)$ 的奇（射影）酉群时，在惰性素数处该空间的维数为 $\begin{pmatrix}\frac{n+1}{2}\\ i\end{pmatrix}$。

对代数闭域上维数为 $d$ 的（未必不可约）代数簇 $Z$，以 $H^{\on{BM}}_{2d}(Z)(-d)$ 记最高次 Borel--Moore 同调的 $(-d)$-Tate 扭，它是以 $Z$ 的不可约分支为基的向量空间。今设 $X$ 是 $q$ 元有限域 $k$ 上维数为 $d+r$ 的光滑簇，$Z \subseteq X_{\overline k}$ 是维数为 $d$ 的（未必不可约）射影子簇。有圈类映射
\[
\cl: H^{\on{BM}}_{2d}(Z)(-d)\to \bigcup_{j\geq 1}H_{c}^{2d}(X_{\overline k}, \bQ_\ell(d))^{\sigma_q^j}=:T^d_\ell(X).
\]

\begin{thm}
\label{T:main theorem}
记 $d_{\mu}=\dim X=2d$，$r=\dim V_\mu^{\Tate}$。
\begin{enumerate}
\item\label{T:main theorem-1} $\overline{\scrS}_{K,\overline{k_v}}$ 的基本 Newton 层 $\overline{\mathscr S}_b$ 是纯维数 $d$ 的，特别地 $d$ 总是整数。此外有 $H_{K,\bQ_\ell}$-等变同构
\[
H^{\on{BM}}_{2d}(\overline{\mathscr S}_{b})(-d)\cong C(G'(\bQ)\backslash G'(\bA_f)/K,\bQ_\ell)^{\oplus r},
\]
其中 $G'$ 是 $G$ 的有限平凡化内形式，且 $G'_\bR$ 紧。
\item\label{T:main theorem-2} 设 $\pi_f$ 是 $H_{K,\Ql}$ 的不可约模，令
\[
H^{\on{BM}}_{2d}(\overline{\mathscr S}_{b})[\pi_f]=\Hom_{H_{K,\Ql}}(\pi_f,H^{\on{BM}}_{2d}(\overline{\mathscr S}_{b})(-d)_{\Ql})\otimes \pi_f
\]
为 $\pi_f$-等型分支。则圈类映射
\[
\cl:H^{\on{BM}}_{2d}(\overline{\mathscr S}_{b})(-d)\to T^d_\ell(\overline{\mathscr S}_{K})
\]
限制在 $H^{\on{BM}}_{2d}(\overline{\mathscr S}_{b})[\pi_f]$ 上单射，只要 $\pi_{f,p}$（即 $\pi_f$ 在 $p$ 处的分支）的 Satake 参数是 $V_\mu$-一般的。
\item\label{T:main theorem-3} 设 $\Sh_K(G)$ 是（本质上的）四元数志村簇或 Kottwitz 算术簇。若 $\pi_{f,p}$ 的 Satake 参数是强 $V_\mu$-一般的，则圈类映射的 $\pi_f$-等型分支满射到 $T^d_\ell(\overline{\scrS_K})[\pi_f]$。特别地，对这些 $\pi_f$，Tate 猜想成立。
\end{enumerate}
\end{thm}

\begin{rmk}
\begin{enumerate}
\item 对 $\hat G$ 的表示 $V$，“$V$-一般”与“强 $V$-一般” Satake 参数的定义见 \cite[Definition 1.4.2]{Xiao-Zhu-Cycle}。$\CG|_{d=(p,\sigma_p)}$ 中的正则半单元总是 $V$-一般的，但反之不然。见 \cite[Remark 1.4.3]{Xiao-Zhu-Cycle}。
\item 该定理的一些特殊情形最初在 \cite{Helm-Tian-Xiao,Tian-Xiao-Tate} 中证明。
\end{enumerate}
\end{rmk}

该定理的证明依赖几个不同的要素。借助志村簇基本部分的 Rapoport--Zink 单值化，第 \eqref{T:main theorem-1} 部分可以化归为关于某些仿射 Deligne--Lusztig 簇的不可约分支的问题，后者在 \cite[\S 3]{Xiao-Zhu-Cycle} 中研究过。最难的是第 \eqref{T:main theorem-2} 部分，我们的证明是计算 $\overline\scrS_b$ 中所有 $d$ 维圈之间的相交数。这些数可以编码为一个元素在 $H^{\cl}_{K_p}$ 中的 $r\times r$ 矩阵。一般说来，直接具体地算出这个矩阵似乎没有指望。

然而，这个矩阵可以理解为 $\Coh(\Loc_{\CG,p}^{\ur})$ 中某些态射的复合。具体地说，先把 $G'(\bQ)\backslash G'(\bA)/K$ 实现为志村集，其志村余特征为 $\mu'=0$。然后利用定理 \ref{introT: spectral action}（以及 Satake 同构 \eqref{E: derived Satake}），该矩阵可计算为
\[
\begin{aligned}
&\Hom_{\Coh(\Loc_{\CG,p}^{\ur})}(\mO,\widetilde{V_{\mu}})\otimes  \Hom_{\Coh(\Loc_{\CG,p}^{\ur})}(\widetilde{V_{\mu}},\mO)\\
&\qquad\to  \Hom_{\Coh(\Loc_{\CG,p}^{\ur})}(\mO,\mO)\cong H^{\cl}_{K_p}\otimes\bQ_\ell.
\end{aligned}
\]
接着需要判定这一配对何时非退化，这本身是表示论中一个有意思的问题，其解决依赖对向量值函数的 Chevalley 限制映射的研究。该矩阵的行列式在 \cite{Xiao-Zhu-Invariants} 中算出。最后，第 \eqref{T:main theorem-3} 部分通过比较两个迹公式（关于 $G$ 的 Lefschetz 迹公式与关于 $G'$ 的 Arthur--Selberg 迹公式）得到证明。

\begin{ex}\label{Tate cycle on U(1,2n)}
设 $G=\on{U}(1,2r)$ 是二次虚扩张 $E/\bQ$ 所关联的 $(2r+1)$ 元酉群\footnote{它不是伴随群，故本例与我们的假设不尽相符；但这样讨论起来更方便，计算本质上也一样。}，其在无穷处的符号为 $(1,2r)$。它带有标准的志村数据，从而（在固定水平 $K\subset G(\bA_f)$ 之后）给出志村簇。特别地，若 $r=1$，这就是（本质上的）Picard 模曲面。设 $p$ 是在 $E$ 中惰性的素数，且 $K_p$ 超特殊。此时 $\overline{\scrS}_b$ 是某些 Deligne--Lusztig 簇之并，由 $G'(\bQ)\backslash G'(\bA_f)/K$ 参数化，其中 $G'=\on{U}(0,2r+1)$ 在所有有限位置同构于 $G$。这些圈在 $\overline{\scrS}_b$ 内部的相交模式（本质上）已由 \cite{Vollaard-Wedhorn-GU(n1)} 给出，但圈之间的相交数要难算得多。事实上，除把定理 \ref{introT: spectral action} 应用于此情形之外，我们并不知道对一般的 $r$ 如何直接计算它们。（$r=1$ 的情形可以直截了当地处理。）

此时 $\hat G=\GL_{2r+1}$，$\sigma_p$ 在其上按 $A\mapsto J(A^T)^{-1}J$ 作用，其中 $J$ 是反对角线上元素全为 $1$ 的反对角矩阵。表示 $V_\mu$ 是 $\GL_{2r+1}$ 的标准表示。可以验证 $\dim V_\mu^{\Tate}=1$（这与上面所说 $\overline{\scrS}_b$ 的不可约分支由 $G'(\bQ)\backslash G'(\bA_f)/K$ 参数化一致）。按惯例把 $\hat G$ 的权格等同于 $\bZ^{2r+1}$。此时 $\Hom_{\Coh(\Loc_{\CG,p}^{\ur})}(\mO,\widetilde{V_{\mu}})$ 是 $\Hom_{\Coh(\Loc_{\CG,p}^{\ur})}(\mO,\mO)=H_{K_p}^{\cl}\otimes\bQ_\ell$ 上的自由秩一 module。于是生成元 $\mathbf a_{\on{in}}$ 诱导出 $H_{K,\bQ_\ell}$-等变同态
\[S(\mathbf a_{\on{in}}): C(G'(\bQ)\backslash G'(\bA_f)/K)\to H_c^{2r}(\overline{\scrS}_{K,\overline{k_v}},\bQ_\ell(r)),\]
它实现 $\overline\scrS_b$ 的圈类映射（差一个倍数）。模 $\Hom_{\Coh(\Loc_{\CG,p}^{\ur})}(\widetilde{V_{\mu}},\mO)$ 同样是 $H_{K_p,\bQ_\ell}$ 上的自由秩一模。对所选生成元 $\mathbf a_{\on{out}}$，复合
\[S(\mathbf a_{\on{out}})\circ S(\mathbf a_{\on{in}})=S(\mathbf a_{\on{out}}\circ\mathbf a_{\on{in}})\]
从 $\overline\scrS_b$ 的不可约分支计算出这些圈的相交矩阵。

元素 $h:=\mathbf a_{\on{out}}\circ\mathbf a_{\on{in}} \in H_{K_p,\bQ_\ell}$ 已在 \cite[Example 6.4.2]{Xiao-Zhu-Invariants} 中显式计算（差显然的修正，并借助 Satake 同构 \eqref{Sat isom classical}）。具体地，
\begin{equation}
\label{E:Hecke operator in Tpj}
h=p^{r(r+1)}\sum_{i=0}^{r} (-1)^i(2i+1)p^{(i-r)(r+i+1)}\sum_{j=0}^{r-i} \begin{bmatrix} 2r+1-2j \\ r-i- j\end{bmatrix}_{t=-p}T_{p,j}.
\end{equation}
这里 $T_{p,j}=1_{K_p\la_j(p)K_p}$，其中 $\la_i=(1^i,0^{2r-2i+1},(-1)^i)$；而 $\begin{bmatrix} n \\ m\end{bmatrix}_t$ 是二项式系数的 $t$-类比，由
\[
[0]_t=1,\quad [n]_t=\frac{t^n-1}{t-1},\quad [n]_t!=[n]_t[n-1]_t\cdots[1]_t,\quad \begin{bmatrix} n \\ m\end{bmatrix}_t=\frac{[n]_t!}{[n-m]_t![m]_t!}
\]
给出。换言之，此时 $\overline\scrS_b$ 中圈的相交矩阵由 Hecke 算子 \eqref{E:Hecke operator in Tpj} 算出。

这一计算有一个有意思的推论，即对有限 Deligne--Lusztig 簇的相交理论的一个结论，而我们并不知道它的直接证明。设 $W$ 是 $\bF_{p^2}$ 上 $(2r+1)$ 维非退化 Hermite 空间。考虑 $r$ 维 Deligne--Lusztig 簇
\[
\mathrm{DL}_{r}: = \big\{ H \subset W \textrm{ 维数为 } r\;\big|\; H \subseteq (H^{(p)})^\perp \big\},
\]
其中 $H^{(p)}$ 是 $H$ 沿 Frobenius 的拉回。以 $\mH$ 记相应的秩 $r$ 万有子丛。令 $\mE = \mH^{(p)} \otimes \big( (\mH^{(p)})^\perp / \mH \big)$。则
\begin{equation}\label{E: chern class on DL}
\int_{\mathrm{DL}_r} c_r( \mE)  = \sum_{i = 0}^r (-1)^i(2i+1) p^{i^2+i} \begin{bmatrix} 2r+1 \\ r-i\end{bmatrix}_{t=-p}.
\end{equation}
\end{ex}

% ============================================================
%  ch03d  §3.4 Rankin--Selberg 动机的 Beilinson--Bloch--Kato 猜想
%  原文对应：ICM-arXiv.tex 1247--1321 行
%  编号：定理 3.4.1 (RS-Sym-Elliptic)；定理 3.4.2 (BBK for RS)
% ============================================================

\subsection{Rankin--Selberg 动机的 Beilinson--Bloch--Kato 猜想}\label{S: BBK}

设 $M$ 是数域 $F$ 上权为 $-1$ 的有理纯 Chow 动机。Beilinson--Bloch--Kato 猜想是 Birch--Swinnerton-Dyer 猜想的深远推广，它预言附于 $M$ 的若干代数、解析与上同调不变量之间的深刻联系：
\begin{itemize}
\item $M$ 的同调平凡圈的有理 Chow 群 $\CH(M)^0$；
\item $M$ 的 $L$-函数 $L(s,M)$；
\item $M$ 的 $\ell$-进实现 $H_{\ell}(M)$ 的 Bloch--Kato Selmer 群 $H^1_f(F, H_{\ell}(M))$。
\end{itemize}
Beilinson--Bloch 猜想预言
\[
\dim_{\bQ}\CH(M)^0 = \on{ord}_{s=0}L(s,M),
\]
即 $\CH(M)^0$ 的维数等于 $L$-函数在中心点的消没阶；而 Bloch--Kato 猜想预言
\[
\on{ord}_{s=0}L(s,M)= \dim_{\bQ_\ell} H^1_f(F, H_{\ell}(M)).
\]
此外，所谓的 $\ell$-进 Abel--Jacobi 映射
\[
\on{AJ}_\ell: \on{CH}(M)^0\otimes \bQ_\ell\to H^1_f(F, H_{\ell}(M))
\]
应当是同构。

该猜想目前看来完全无从下手。例如，对一般的动机，$L$-函数是否能亚纯延拓到整个复平面（从而 $L(s,M)$ 在 $s=0$ 处的消没阶才有意义）都还远未解决。（这一点会由数域上朗兰兹对应的 Galois 到自守方向推出。）尽管如此，已有许多工作在各种特殊情形下检验该猜想，多数是针对低秩的动机 $M$。在 \cite{LTXZZ-BBK} 中，我们对 Rankin--Selberg 动机验证了上述猜想的某些情形，而这类动机构成一列秩任意大的动机。

以下设 $F/F^+$ 是（非平凡的）CM 扩张，$F^+$ 全实。

\begin{thm}\label{RS-Sym-Elliptic}
设 $A_1,A_2$ 是 $F^+$ 上的两条椭圆曲线，并假设
\begin{enumerate}
\item\label{RS-Sym-Elliptic-1} $\End_{\overline F}A_i=\bZ$ 且 $\Hom_{\overline F}(A_1,A_2)=0$；
\item\label{RS-Sym-Elliptic-2} $\Sym^{n-1}A_1$ 与 $\Sym^nA_2$ 是模的；
\item\label{RS-Sym-Elliptic-3} 若 $n\geq 3$，则 $F^+\neq \bQ$。
\end{enumerate}
在这些假设下，若 $L(n, \Sym^{n-1}A_1\times \Sym^nA_2)\neq 0$，则对几乎所有的 $\ell$ 有
\[
\dim_{\bQ_\ell} H^1_f(F,  \Sym^{n-1}V_\ell(A_1)\otimes\Sym^nV_\ell(A_2)(1-n))=0.
\]
这里 $V_\ell(A_i)$ 照例记 $A_i$ 的有理 Tate 模。
\end{thm}

该定理其实是关于某些 Rankin--Selberg 自守表示所关联的 Galois 表示的 Bloch--Kato Selmer 群的一个更一般结果的特例，下面就来讨论它。

回忆一下，对 $\GL_{n}(\bA_F)$ 的一个不可约正则代数共轭自对偶尖（RACSDC）自守表示 $\Pi$，可以附上一族相容的 Galois 表示 $\rho_{\Pi,\la}:\Ga_F\to \GL_n(E_\la)$，其中 $E\subset \bC$ 是足够大的数域，$\la$ 是 $E$ 的一个素位（见 \cite{Chenevier-Harris-Galois-Rep}）。这样的 $E$ 称为 $\Pi$ 的强系数域；在下面考虑的情形中，它可以取为 $\Pi$ 的 Hecke 特征值生成的数域。

\begin{thm}\label{BBK for RS}
若 $n\geq 3$，设 $F^+\neq\bQ$。设 $\Pi_n$（相应地 $\Pi_{n+1}$）是 $\GL_{n}(\bA_F)$（相应地 $\GL_{n+1}(\bA_F)$）的 RACSDC 自守表示，无穷小特征平凡。设 $E\subseteq\bC$ 同时是 $\Pi_n$ 与 $\Pi_{n+1}$ 的强系数域。设 $\lambda$ 是 $E$ 的关于 $\Pi:=\Pi_0\times \Pi_1$ 的可容许素位\zy{原文此处把 $\Pi_n\times\Pi_{n+1}$ 写作 $\Pi_0\times\Pi_1$，按上下文应为前者；见勘误表。}。令 $\rho_{\Pi,\la}:=\rho_{\Pi_n,\la}\otimes_{E_\la} \rho_{\Pi_{n+1},\la}$。
\begin{enumerate}
\item\label{BBK for RS-1} 若 Rankin--Selberg $L$-值 $L(\frac{1}{2},\Pi)\neq 0$\footnote{这里用的是 $L$-函数的自守归一化。}，则 $H^1_f(F,\rho_{\Pi,\lambda}(n))=0$。
\item\label{BBK for RS-2} 若某个元素 $\Delta_\la\in H^1_f(F,\rho_{\Pi,\lambda}(n))$（下面会解释）非零，则 $H^1_f(F,\rho_{\Pi,\lambda}(n))$ 作为 $E_\la$-向量空间由 $\Delta_\la$ 生成。
\end{enumerate}
\end{thm}

定理中出现的可容许素位的概念包含一长串假设，其中一些相当技术性。本质上，它保证 Galois 表示 $\rho_{\Pi,\lambda}$ 有定义良好的 $\mO_{E,\la}$-格（下文仍记作 $\rho_{\Pi,\lambda}$），且模 $\la$ 约化表示足够大，并含某些特定的元素。（这也与定理 \ref{T:main theorem} 中 $V$-一般的概念有关。）所幸在一些有利情形下，可以证明除有限多个素位外都是可容许的。例如定理 \ref{RS-Sym-Elliptic} 所考虑的情形就是如此。纯自守框架下的另一个情形见 \cite[Thm. 1.1.7]{LTXZZ-BBK}。

该定理的证明用到几个不同的要素。情形 \eqref{BBK for RS-1} 的第一步，是借助整体 Gan--Gross--Prasad（GGP）猜想，把解析条件 $L(\frac{1}{2},\Pi)\neq 0$ 转化为更代数的条件。也就是说，GGP 猜想预言此时存在 $F$ 上的一对 Hermite 空间 $(V_n,V_{n+1})$，在 $\infty$ 处全正定，且 $V_{n+1}=V_n\oplus Fv$，其中 $(v,v)=1$；还存在酉群之积 $G=U(V_n)\times U(V_{n+1})$ 的一个温度尖自守表示 $\pi=\pi_n\times \pi_{n+1}$，使得周期积分
\[
[\Delta_H]: C^*_c(\Sh(G),E)[\pi]\to E
\]
不消失。这里 $H:=U(V_n)$ 对角地嵌入 $G$，从而诱导出相应志村簇（实际上是志村集）的嵌入 $\Delta_H: \Sh(H)\hookrightarrow \Sh(G)$，水平结构取适当（这里及以后都从记号中略去水平结构）。以 $[\Delta_H]$ 记 $\Sh(G)$ 中由 $\Sh(H)$ 给出的同调类，以 $C^*_c(\Sh(G),E)[\pi]$ 记 $\Sh(G)$ 的上同调（即函数空间）的 $\pi$-等型分支。

该猜想最初在 \cite{Zhang-GGP} 中，在某些局部限制下得到证明，而那些限制对算术应用来说过强。我们最近的工作借助迹公式研究中的一些新技术，把这些限制全部去掉了（\cite{BLZZ-GGP}）。

接下来的策略是：对每个 $m\geq 1$ 构造（无穷多个）上同调类 $\{\Theta_m^p\}_p\subset H^1(F, (\rho_{\Pi,\la}/\la^m)^*(1))$，其中 $p$ 是适当选取的素数，$(-)^*(1)$ 记 Pontryagin 对偶再扭以分圆特征；使得 $\Theta_m^p$ 与 Galois 表示 $\rho_{\Pi,\la}/\la^m$ 的 Selmer 类在 $p$ 处的局部 Tate 配对与上述周期积分相关。然后就可以用 Kolyvagin 型论证（在 \cite{Liu-HZ-Cycle, LTXZZ-BBK} 中作了放大），以 $\{\Theta_m^p\}$ 作为 Selmer 群的零化子，从而得出结论。

$\Theta_m^p$ 的构造用到志村簇的对角嵌入
\[
\Delta_{H'}: \Sh(H')\hookrightarrow \Sh(G'),
\]
其中 $H'\hookrightarrow G'$ 是 $H\subset G$ 的 prime-to-$p$ 平凡化（扩张纯）内形式（见定义 \ref{D: prime-to-p trivial}）。这些志村簇在 $p$ 处有抛孔水平结构，利用局部模型理论（\S\ref{S: local model}）可以证明它们的整模型在 $p$ 处是多重半稳定的，并算出其 mod $p$ 纤维上的邻近圈层。

综合运用许多要素——包括对 $\Sh(G')$ 在 $\overline{F}$ 上（整）上同调的理解、例 \ref{Tate cycle on U(1,2n)} 中的计算（特别是 \eqref{E:Hecke operator in Tpj} 与 \eqref{E: chern class on DL}）、以及 Taylor--Wiles 拼接方法 \cite{LTXZZ-Rigidity}——可以证明 $(\rho_{\Pi,\la}/\la^m)^*(1)$ 确实出现在 $\Sh(G')$ 的上同调中（即所谓的 $\Pi$ 的算术提高水平），并且对角圈 $\Delta_{H'}$ 在局部化到 $(\rho_{\Pi,\la}/\la^m)^*(1)$ 时确实给出所要的类 $\Theta_m^p$。需要说明的是，这与 \S\ref{S: Local Langlands} 与 \S\ref{geometric JL} 中的猜想是相容的：在 $\Loc_{\CG,p}\otimes \mO_E/\la^m$ 上，对应于 $c\on{-ind}_{K_p}^{G(\bQ_p)}(\mO_E/\la^m)$ 与 $c\on{-ind}_{K'_p}^{G'(\bQ_p)}(\mO_E/\la^m)$ 的凝聚层正应当以这种方式相关联。

情形 \eqref{BBK for RS-2} 中出现的类 $\Delta_\la$ 也可以解释。此时我们从一个志村簇的嵌入 $\Delta_H: \Sh(H)\hookrightarrow\Sh(G)$ 出发，其中 $G$ 是酉群之积，使得 $\Pi$ 下降为出现在 $\Sh_G$ 中间维上同调中的温度尖自守表示 $\pi$。于是圈 $\Delta_H$ 的 $\pi$-等型分支同调于零，我们令 $\Delta_\la=\on{AJ}_\la(\Delta_H[\pi])$。证明情形 \eqref{BBK for RS-2} 的策略，则是通过与前面类似的论证把它化归为情形 \eqref{BBK for RS-1}。

% ============================================================
%  ch04  参考文献（原文 1345--1553 行；条目保留原文，仅统一排版）
% ============================================================

\clearpage
\phantomsection
\addcontentsline{toc}{section}{参考文献}

\begin{thebibliography}{99}

\bibitem{AGLR}J.~Ansch\"utz and I.~Gleason and J.~Louren\c{c}o and T.~Richarz, On the $p$-adic theory of local models. ArXiv: 2201.01234.

\bibitem{Arinkin-Gaitsgory-Singular-Support}D.~Arinkin and D.~Gaitsgory, Singular support of coherent sheaves, and the geometric Langlands conjecture. \emph{Selecta Math.} \textbf{21} (2015), 1--199.

\bibitem{AGKRRV}D.~Arinkin and D.~Gaitsgory and D.~Kazhdan and S.~Raskin and N.~Rozenblyum and Y.~Varshavsky, The stack of local systems with restricted variation and geometric Langlands theory with nilpotent singular support. ArXiv:2010.01906.

\bibitem{Beauville-Laszlo}A.~Beauville and Y.~Laszlo, Conformal blocks and generalized theta functions. \emph{Comm. Math. Phys.} \textbf{164} (1994), no.~2, 385--419.

\bibitem{Beilinson-Drinfeld-Hitchin}A.~Beilinson and V.~Drinfeld, Quantization of Hitchin's integrable system and Hecke eigensheaves. \href{www.math.uchicago.edu/~mitya/langlands}.

\bibitem{BCHN}D.~Ben-Zvi and H.~Chen and D.~Helm and D.~Nadler, Coherent Springer theory and the categorical Deligne-Langlands correspondence. ArXiv:2010.02321.

\bibitem{Ben-Zvi-Nadler-Betti-Geometric-Langlands}D.~Ben-Zvi and D.~Nadler, Betti geometric Langlands. In \emph{Algebraic Geometry: Salt Lake City 2015}, edited by T.~de Fernex et al., pp. 3--41, Proc. Symp. Pure Math. 97.2, Amer. Math. Soc., Providence, RI, 2018.

\bibitem{BLZZ-GGP} R.~Beuzart-Plessis and Y.~Liu and W.~Zhang and X.~Zhu. Isolation of cuspidal spectrum, with application to the Gan-Gross-Prasad conjecture. \emph{Ann. Math.} \textbf{194} (2021), no.~2, 519--584.

\bibitem{Bezrukavnikov-Affine-Hecke}R.~Bezrukavnikov, On two geometric realizations of an affine Hecke algebra. \emph{Publ. Math. IH\'ES} \textbf{123}, (2016) 1--67.

\bibitem{Bezrukavnikov-Braverman-Geometric-Langlands-char-p}R.~Bezrukavnikov and A.~Braverman, Geometric Langlands
conjecture in characteristic p: The $GL_n$ case. \emph{Pure Appl. Math. Q.} \textbf{3} (2007), no.~1, Special Issue: In honor of Robert D. MacPherson, Part 3, 153--179.

\bibitem{Bezrukavnikov-Finkelberg-Derived-Satake}R.~Bezrukavnikov and M.~Finkelberg, Equivariant Satake category and Kostant-Whittaker reduction. \emph{Mosc. Math. J.} \textbf{8} (2008), no.~1, 39--72.

\bibitem{Bezrukavnikov-Travkin-Quantization}R.~Bezrukavnikov and R.~Travkin and T.-H.~Chen and X.~Zhu, Quantization of Hitchin integrable system via positive characteristic. ArXiv:1603.01327.

\bibitem{Bhatt-Scholze-Witt}B.~Bhatt and P.~Scholze, Projectivity of Witt vector affine Grassmannians. \emph{Invent. Math.} \textbf{209} (2017), no.~2, 329--423.

\bibitem{Chen-Zhu-Nonabelian-Hodge}T.-H.~Chen and X.~Zhu, Non-abelian Hodge theory for curves in characteristic $p$. \emph{Geometric and Functional Analysis} \textbf{25} (2015), 1706--1733.

\bibitem{Chen-Zhu-Geometric-Langlands-char-p}T.-H.~Chen and X.~Zhu, Geometric Langlands in prime characteristic. \emph{Comp. Math.} \textbf{153} (2017), 395--452.

\bibitem{Chenevier-Harris-Galois-Rep}G.~Chenevier and M.~Harris, Construction of automorphic Galois representations, II, \emph{Camb. J. Math.} \textbf{1} (2013), no.~1, 53--73.

\bibitem{DHKM}J.F.~Dat and D.~Helm and R.~Kurinczuk and G.~Moss, Moduli of Langlands parameters. ArXiv: 2009.06708.

\bibitem{DLYZ-Endoscopy}G.~Dhillon and Y.~Li and Z~Yun and X.~Zhu, in preparation.

\bibitem{Donagi-Pantev-Hitchin}R.~Donagi and T.~Pantev, Langlands duality for Hitchin systems. \emph{Invent. Math.} \textbf{189} (2012), 653--735.

\bibitem{Drinfeld-GL2}V. G.~Drinfeld, Two-dimensional $\ell$-adic representations of the fundamental group of a curve over a finite field and automorphic forms on $\GL(2)$, \emph{Amer. J. Math.} \textbf{105} (1983), 85--114.

\bibitem{Emerton-Gee-Zhu}M.~Emerton and T.~Gee and X.~Zhu, in preparation.

\bibitem{Emerton-Zhu}M.~Emerton and X.~Zhu, in preparation.

\bibitem{Faltings-Loop}G.~Faltings, Algebraic loop groups and moduli spaces of bundles. \emph{J. Eur. Math. Soc.} \textbf{5} (2003), no.~1, 41--68.

\bibitem{Faltings-Chai-Abelian}G.~Faltings and C.L.~Chai, \emph{Degeneration of abelian varieties}. volume 22 of Results in Mathematics and Related Areas (3), Springer-Verlag, Berlin, 1990.

\bibitem{Fargues-Scholze-Geometrization}L.~Fargues and P.~Scholze, Geometrization of the local Langlands correspondence. ArXiv:2102.13459.

\bibitem{FGV-Geometric-Langlands-GLn}E.~Frenkel and D.~Gaitsgory and K.~Vilonen, On the geometric Langlands conjecture. \emph{J. Amer. Math. Soc.} \textbf{15} (2001), 367--417.

\bibitem{Gaitsgory-Central-Sheaf}D. Gaitsgory, Construction of central elements in the affine Hecke algebra via nearby cycles. \emph{Invent. math.} \textbf{144} (2001), 253--280.

\bibitem{Gaitsgory-deJong}D. Gaitsgory, On De Jong's conjecture. \emph{Isr. J. Math.} \textbf{157} (2007), 155--191.

\bibitem{Haines-Test-Function}T.~Haines, The stable Bernstein center and test functions for Shimura varieties. In \emph{Automorphic forms and Galois representations. Vol. 2}, edited by
F.~Diamond, P.L.~Kassaei and M.~Kim, pp.118--186, London Math. Soc. Lecture Note Ser. 415, Cambridge Univ. Press, Cambridge, 2014.

\bibitem{Harris-Taylor-GLn}M.~Harris and R.~Taylor, \emph{The Geometry and Cohomology of Some Simple Shimura Varieties}.  Ann. Math. Studies 151, Princeton University Press, 2001.

\bibitem{HNY-Kloosterman}J.~Heinloth and B.C.~Ng\^o and Z.~Yun, Kloosterman sheaves for reductive groups. \emph{Ann. Math.} \textbf{177} (2013), 241--310.

\bibitem{Hellmann-Iwahori}E. Hellmann, On the derived category of the Iwahori-Hecke algebra. ArXiv:2006.03013.

\bibitem{Helm-Tian-Xiao}D.~Helm and Y.~Tian and L.~Xiao, Tate cycles on some unitary Shimura varieties mod $p$. \emph{Algebraic Number Theory} \textbf{11} (2017), 2213--2288.

\bibitem{HRS-Kunneth}T.~Hemo and T.~Richarz and J.~Scholbach, Constructible sheaves on schemes and a categorical K\"unneth formula. ArXiv:2012.02853.

\bibitem{Hemo-Zhu-Unipotent}T.~Hemo and X.~Zhu, Unipotent categorical local Langlands correspondence. Preprint.

\bibitem{Kazhdan-Lusztig-Iwahori}D.~Kazhdan and G.~Lusztig, Proof of the Deligne-Langlands conjecture for Hecke algebras.\emph{Invent. Math.} \textbf{87} (1987), 153--215.

\bibitem{Kisin-Integral-Model}M.~Kisin, Integral models for Shimura varieties of abelian type. \emph{J. Amer. Math. Soc.} \textbf{23} (2010), 967--1012.

\bibitem{Lafforgue-GLn}L.~Lafforgue, Chtoucas de Drinfeld et correspondance de Langlands. \emph{Invent. math.} \textbf{147} (2002), 1--241.

\bibitem{Lafforgue-Automorphic-to-Galois}V.~Lafforgue, Chtoucas pour les groupes r\'eductifs et param\`etrisation de Langlands globale. \emph{J. Amer. Math. Soc.} \textbf{31} (2018), 719--891.

\bibitem{Lafforgue-Zhu-Elliptic}V.~Lafforgue and X.~Zhu, D\'ecomposition au-dessus des param\`etres de Langlands elliptiques. ArXiv: 1811.07976.

\bibitem{Langlands-to-Weil}R.~Langlands, Letter to Andr\'e Weil. \href{publications.ias.edu/rpl/paper/43}.

\bibitem{Laumon-Geometric-Langlands}G.~Laumon, Correspondance de Langlands g\'eom\'etrique pour les corps de fonctions.\emph{Duke Math. J.} \textbf{54} (1987), 309--359.

\bibitem{LRS-GLn}G.~Laumon and M.~Rapoport and U.~Stuhler, $D$-elliptic sheaves and the Langlands correspondence. \emph{Invent. Math.} \textbf{113} (1993), 217--338.

\bibitem{LLLM}D.~Le and B.V.~Le Hung and B.~Levin and S. Morra, Local models for Galois deformation rings and applications. ArXiv:2007.05398.

\bibitem{Lee-ES-Relation}S.Y.~Lee, Eichler-Shimura Relations for Shimura Varieties of Hodge Type. ArXiv:2006.11745,

\bibitem{Levin-Weil-Restriction}B.~Levin, Local models for Weil-restricted groups. \emph{Comp. Math.}, \textbf{152} (2016), 2563--2601.

\bibitem{Liu-HZ-Cycle}Y.~Liu, Hirzebruch-Zagier cycles and twisted triple product Selmer groups. \emph{Invent. Math.} \textbf{205} (2016) no.~3, 693--780.

\bibitem{LTXZZ-Rigidity}Y.~Liu and Y.~Tian and L.~Xiao and W.~Zhang and X.~Zhu, Deformation of rigid conjugate self-dual Galois representations. ArXiv:2108.06998.

\bibitem{LTXZZ-BBK}Y.~Liu and Y.~Tian and L.~Xiao and W.~Zhang and X.~Zhu, On the Beilinson-Bloch-Kato conjecture for Rankin-Selberg motives. \emph{Invent. Math} (2022).

\bibitem{Lourenco}J.~Louren\c{c}o, Grassmanniennes affines tordues sur les entiers. ArXiv:1912.11918.

\bibitem{Lusztig-q-analog}G.~Lusztig, Singularities, character formulas, and a $q$-analog of weight multiplicities. In \emph{Analysis and topology on singular spaces, II, III (Luminy, 1981)}, pp. 208--229, Ast\'erisque, 101-102, Soc. Math. France, Paris, 1983.

\bibitem{Lusztig-Unipotent}G.~Lusztig, Classification of unipotent representations of simple $p$-adic groups. \emph{Int. Math. Res. Not.} (1995), 517--589.

\bibitem{Mirkovic-Vilonen-GeoSat}I.~Mirkovi\'c and K.~Vilonen, Geometric Langlands duality and representations of algebraic groups over commutative rings. \emph{Ann. Math.} \textbf{166} (2007), 95--143.

\bibitem{Nekovar-Scholl-Plectic}J.~Nekov\'a\v{r} and A.J.~Scholl, Introduction to plectic cohomology. In \emph{Advances in the theory of automorphic forms and their L-functions}, Contemp. Math. Vol. 664, pp 321--337. Amer. Math. Soc., Providence, RI, 2016.

\bibitem{Ngo-Fundamental-Lemma}B.C.~Ng\^o, Le lemme fondamental pour les alg\`ebres de Lie. \emph{Publ. Math. l'IH\'ES} \textbf{111} (2010), 1--169.

\bibitem{Pappas-Rapoport-Ramified-Unitary}G.~Pappas and M.~Rapoport, Local models in the ramified case. III. Unitary groups.  \emph{J. Inst. Math. Jussieu} \textbf{8} (2009), 507--564.

\bibitem{PRS-Survey}G.~Pappas and M.~Rapoport and B.D.~Smithling, Local models of Shimura varieties, I. Geometry and combinatorics. In \emph{Handbook of Moduli, Vol. III}, pp. 135--217, International Press, Somerville, MA, 2013.

\bibitem{Pappas-Zhu-Local-Model}G.~Pappas and X.~Zhu, Local models of of Shimura varieties and a conjecture of Kottwitz.\emph{Invent. math.} \textbf{194} (2013), 147--254.

\bibitem{Scholze-Berkeley}P.~Scholze and J.~Weinstein, \emph{Berkeley lectures on $p$-adic geometry}. Ann. Math. Studies 207, Princeton University Press, 2020.

\bibitem{Tian-Xiao-Tate}
Y.~Tian and L.~Xiao, Tate cycles on quaternionic Shimura varieties over finite fields. \emph{Duke Math. J.} \textbf{168} (2019), 1551--1639.

\bibitem{Vollaard-Wedhorn-GU(n1)}I.~Vollaard and T.~Wedhorn, The supersingular locus of the Shimura variety of $GU(1,n-1)$, II. \emph{Invent. Math.} \textbf{184} (2011), 591--627.

\bibitem{Wiles-Fermat}A.~Wiles, Modular elliptic curves and Fermat's Last Theorem. \emph{Ann. Math.} \textbf{141} (1995), 443--551.

\bibitem{Wu-S=T}Z.~Wu, $S=T$ for Shimura Varieties and $p$-adic Shtukas. ArXiv:2110.10350.

\bibitem{Xiao-Zhu-Cycle}L.~Xiao and X.~Zhu, Cycles on Shimura varieties via geometric Satake. ArXiv:1707.05700.

\bibitem{Xiao-Zhu-Invariants}L.~Xiao and X.~Zhu, On vector-valued twisted conjugate invariant functions on a group. In \emph{Representations of Reductive Groups}, pp. 361--426, Proc. Symp. Pure Math. 101, Amer. Math. Soc., Providence, RI, 2019.

\bibitem{Xu-Zhu-Bessel}D.~Xu and X.~Zhu, Bessel $F$-isocrystals for reductive groups. In \emph{Invent. Math} (2022).

\bibitem{Xue-Finiteness}C.~Xue, Finiteness of cohomology groups of stacks of shtukas as modules over Hecke algebras, and applications. \emph{\'Epijournal de G\'eom\'etrie Alg\'ebrique} \textbf{4} (2020).

\bibitem{Xue-Smoothness}C.~Xue, Smoothness of cohomology sheaves of stacks of shtukas. ArXiv:2012.12833.

\bibitem{Yu-Integral-Satake}J.~Yu, The integral geometric Satake equivalence in mixed characteristic. ArXiv:1903.11132.

\bibitem{Yu-Jacquet-Langlands-Higher-Weights}J.~Yu, A geometric Jacquet-Langlands Transfer for automorphic forms of higher weights. In \emph{Trans. of Amer. Math. Soc.} (2022).

\bibitem{Yun-motive-exceptional-group}Z.~Yun, Motives with exceptional Galois groups and the inverse Galois problem. \emph{Invent. Math.} \textbf{196} (2014), 267--337.

\bibitem{Yun-Zhu-Integral-Homology}Z.~Yun and X.~Zhu, Integral homology of loop groups via Langlands dual groups. \emph{Represent. Theory} \textbf{15}, (2011), 347--369.

\bibitem{Zhang-GGP}W.~Zhang, Fourier transform and the global Gan-Gross-Prasad conjecture for unitary groups. \emph{Ann. Math.} \textbf{180} (2014), 971--1049.

\bibitem{Zhu-S=T}X.~Zhu, $S=T$ for Shimura varieties. Preprint.

\bibitem{Zhu-Coherence-Conjecture}X.~Zhu, On the coherence conjecture of Pappas and Rapoport. \emph{Ann. Math.} \textbf{180} (2014), 1--85.

\bibitem{Zhu-Ramified-Sat}X.~Zhu, The Geometric Satake Correspondence for Ramified Groups. \emph{Ann. Sci. \'Ec. Norm. Sup\'er.}, \textbf{48} (2015),  409--451.

\bibitem{Zhu-CDM}X.~Zhu, Geometric Satake, categorical traces, and arithmetic of Shimura varieties. In \emph{Current Developments in Mathematics} (2016), 145--206.

\bibitem{Zhu-Satake-Mixed}X.~Zhu, Affine Grassmannians and the geometric Satake in mixed characteristic. \emph{Ann. Math.} \textbf{185} (2017), 403--492.

\bibitem{Zhu-FG-vs-HNY}X.~Zhu, Frenkel-Gross' irregular connection and Heinloth-Ng\^o-Yun's are the same. \emph{Selecta Math. (N.S.)} \textbf{23} (2017), no.~1, 245--274.

\bibitem{Zhu-PCMI}X.~Zhu, An introduction to affine Grassmannians and the geometric Satake equivalence. In \emph{IAS/Park City Math. Ser.} \textbf{24} (2017), 59--154.

\bibitem{Zhu-ICCM}X.~Zhu, A survey of local models of Shimura varieties. In \emph{Proceedings of the Sixth International Congress of Chinese Mathematicians}, pp. 247--265, Adv. Lectures in Math. Vol. 36, International Press , Boston, MA, 2017.

\bibitem{Zhu-Conjecture}X.~Zhu, Coherent sheaves on the stack of Langlands parameters. ArXiv:2008.02998.

\bibitem{Zhu-Integral-Satake}X.~Zhu, A note on integral Satake isomorphisms. ArXiv:2005.13056.

\end{thebibliography}


\end{document}
